\chapter{Executable Predictive Memory}\label{sec:memory}

\section{Which memory is being minimized?}
A rank-$r$ linear realization stores $r$ real coordinates. This statement
alone does not determine how many bits an executable approximation needs,
how much immutable code it uses, or how expensive its update and readout
are. The purpose of this section is to separate those quantities and give
exact results in interfaces where they can be compared rigorously.

The equality between finite Hankel rank and minimum linear realization
dimension is classical~\citep{carlyle1971,fliess1974}.
The underlying linear convention is
$f(w)=\alpha^\top A_w\beta$, with row state
$s(w)=\alpha^\top A_w$. Words are read chronologically; geometric
address strings are stored newest first and explicitly reversed for
replay. Logarithms used for bit counts are $\log_2$. Natural logarithms
appearing in pressure or analytic estimates are indicated by $\log$.
We distinguish recurrent dimension, persistent identifier bits,
immutable program size, workspace, numerical output precision and time.
An uncharged clock is not allowed in autonomous finite-state claims.

The operator norm in Section~\ref{sec:hankel} and the completed-word
divergences in Section~\ref{sec:probability} measure different errors
from uniform conditional prediction. We use the shared foundations in
Subsections~\ref{I:U:sec:preliminaries}, \ref{I:S:sec:exact},
and~\ref{I:S:sec:analytic} for the constructions below.

\section{Exact residual state behind a suffix window}
Fix a window length $N$ and write $q_N(w)=\operatorname{suff}_N(w)$,
retaining the whole word when $|w|<N$. Define
\[
 U_N=\operatorname{span}\{s(u)-s(v):q_N(u)=q_N(v)\}.
\]
The space $U_N$ consists of differences the external window fails to
identify. Its transition closure, rather than the ambient Hankel rank,
determines the residual state in the fixed affine interface
\[
 z(wa)=z(w)B_a+g(q_N(w),a),\qquad
 f(w)=z(w)c+h(q_N(w)).
\]
Here the matrices $B_a$ and readout $c$ are fixed, while the window and
forcing tables are supplied externally.

\begin{theorem}[Residual dimension for a fixed affine interface]
For a minimal realization,
\[
 U_N=\sum_{|q|=N}U_0A_q,\qquad
 U_{N+1}=\sum_a U_NA_a\subseteq U_N.
\]
The first plateau persists. The minimum recurrent dimension in the
specified interface is $\dim U_N$. The minimum Loewy length is that
of the actual letter-marked transition algebra acting on $U_N$.
\end{theorem}
Theorems~\ref{III:S:thm:filtration} and \ref{III:S:thm:window} give the exact
construction and lower bound. For the latter, differences within a
window fiber cancel both forcing and direct readout. Equality of all
future outputs then induces a surjective letter-intertwining map onto
the canonical difference space. Choosing its basis attains the bound.

If the homogeneous update and readout may depend on the window, the
correct interface has per-window dimensions $r_q$ and rectangular
$B_{q,a}\in\mathbb K^{r_q\times r_{q'}}$. Its charged physical
quantity is $\max_q r_q$, since only one window is active. Defining
$F_N(q)=\operatorname{span}\{s(u)-s(v):q_N(u)=q_N(v)=q\}$, the
minimum is $\max_q\dim F_N(q)$, with a marked path-image algebra
rather than the fixed-letter algebra. Definition~\ref{III:S:def:fiber-interface}
and Theorem~\ref{III:S:thm:fiber} make this contract explicit. The two
interfaces optimize different quantities.

An invertible neutral letter may preserve every $U_N$, even when
another letter has a nilpotent action. Thus single-letter forgetting does
not certify a memory saving under arbitrary ordered input sequences.
Conditional future laws introduce another distinction: normalization can
preserve differences that decay in the unnormalized linear state.

\section{Persistent bits and exact finite-state examples}
For a declared output error contract, an autonomous machine consists of
$(Q,q_0,(U_a),D)$. The charged persistent identifier length is
$B=\lceil\log_2|Q|\rceil$. Two histories with a future output difference
strictly larger than $2\epsilon$ cannot end in the same state of an
$\epsilon$-accurate machine. This classical packing principle becomes an
exact finite result only after a matching executable construction is given.

For the explicit rank-two family $f_a$ of
Section~\ref{III:M:sec:memory}, the manuscript proves
\[
 B_{f_a}^*(2^{-((a+1)m+2)})=m+2,
 \qquad a,m\in\mathbb Z,\quad a\ge1,\quad m\ge a+2.
\]
This is exact in integer bits. It does not identify the exact minimum
number of states; distinct state counts can have the same ceiling of
their base-two logarithm. The construction packs its clock into the state
identifier. The leading memory coefficient varies with $a$ although
Hankel rank stays two, so rank alone cannot determine that coefficient.

For the unequal-contraction system
\[
 \tau_0(x)=x/4,\qquad \tau_1(x)=x/2+1/2,\qquad h(x)=x,
\]
let $\theta=2^{-\ell}$ and let $\mathcal L_\ell$ be the stopped
newest-first address leaves. Put
$U_\ell=\bigcup_{v\in\mathcal L_\ell}\tau_v([0,1])$.
The all-history contract includes the empty history and uses no external clock.

\begin{theorem}[Fibonacci optimum on the leaf union]
For $\ell\ge1$, initial point $x_0\in U_\ell$, and error at most
$\theta/2$ at every finite history,
\[
 |Q|_{\min}=F_{\ell+2},\qquad
 B_{\min}=\lceil\log_2F_{\ell+2}\rceil.
\]
For every $x_0\in[0,1]$, the same state lower bound holds.
An exact count outside $U_\ell$ is not established here.
\end{theorem}
Theorem~\ref{III:M:rev13:thm:fib-general} proves this statement. Stopped
leaves supply the finite automaton and midpoint outputs. A family of
same-length histories approaching separated leaf endpoints forces the
matching lower bound. The initial-point restriction is necessary for the
stated leaf-machine initialization and is not removed by an asymptotic
dimension calculation.

\section{Rational addresses and irreversible modes}
For rational contractions, a stopped address can be ranked and unranked
without storing an exponential transition table. Exact integer counting
controls the persistent identifier. The proofs charge update arithmetic,
startup and workspace separately; a bit saving does not automatically
change byte-aligned resident memory.

Now let a finite mode obey $r_{t+1}=G(r_t,a_{t+1})$, allowing irreversible
resets, and let the contraction ratio depend on both letter and previous
mode. Storing only the retained letters can lose the information needed
to replay their dynamics. The construction stores a checkpoint
$b_t=r_{t-|v_t|}$ at the beginning of the retained window, together with
its newest-first address $v_t$. Forward replay recovers the applied mode
sequence. The algorithm then truncates using the actual product of
ratios and advances the checkpoint through discarded oldest letters.
It never requires inversion of $G$.

The pressure matrix is
\[
 M(s)_{m,m'}=\sum_{a:G(m,a)=m'}r_{a,m}^{s}.
\]
For the strongly connected mode graph, the positive root
$\rho(M(s^*))=1$ controls the symbolic address count. Theorem~\ref{III:M:rev13:thm:mode}
gives the construction and its conditional optimality statement with
all geometric and future-observability assumptions retained. Symbolic
pressure need not equal geometric dimension under arbitrary overlaps;
Hausdorff dimension by itself does not supply a finite executable with
constant overhead.

Numerical guards are sound only when the certified state enclosure lies
on one side of the decision threshold. If it intersects both sides, the
result is inconclusive. Refinement does not guarantee eventual resolution
at a genuinely ambiguous point. The feedback results certify cache
choices at the time of use and distinguish bounded returned errors from
unconditional discovery of persistent semantic facts.

\section{Execution, transfer approximations and actual coding cost}
An integer-execution certificate propagates roundoff through an explicit
invariant error bound. For transfer operators, the real branches must
preserve $[-1,1]$ and satisfy the stated complex two-radius bounds.
Repeated coefficient truncation then has the finite-horizon guarantee
\[
 \|\mathcal L_\sigma^n1-P_n\|_\infty
 \le n\delta_m K^{n-1},
 \qquad P_{n+1}=\Pi_m\mathcal L_\sigma P_n.
\]
This includes errors accumulated after each truncation; replacing
$P_n$ by a single final truncation is invalid in general.
Theorem~\ref{III:U:thm:oracle} and the rounding theorem state the additional
positivity, normalization and arithmetic hypotheses. Sound finite-horizon
bounds do not imply that a proposed compiler always terminates or that
its total bit complexity is optimal.

For recurrent matrix states, an exact main-state-plus-residual identity
is useful only after the serialized representation is charged. A smaller
residual norm, covariance or rank is not itself a smaller lossless code.
The final application develops reversible integer lifting with an actual
coding contract and separates complete encoded length from execution
workspace and update time. The application includes a controlled example
and states the conditions needed for a measured coding gain.

\section{Open scope and proof organization}
Open problems include exact Fibonacci counts outside the specified leaf
union, general nonconformal memory laws with effective updates, discovery
of protected facts from observations, and compiler completeness with a
sufficient precision rule and full cost bounds. A demonstrated advantage
in complete lossless code length and execution time on independent
recurrent trajectories is also still required for the codec application.

The remaining subsections cover exact residual interfaces,
conditional-future analysis, persistent-bit and address constructions,
mode-aware replay and pressure bounds, numerical execution certificates,
and the recurrent-state coding application. Every theorem is tied to its
own interface and error contract.
\label{III:III:main-end}

\section*{Detailed proofs and supporting results}

\section{Exact memory with an external suffix window}\label{III:S:sec:window}
This section is algebraic and does not assume bounded $H_f$. In a minimal row-state realization over $\SK$, let $q_N(w)=\Ssuff_N(w)$; for $|w|<N$, retain the entire word. The external window uses the original alphabet. Define
\[
 U_N=\SSpan\{s(u)-s(v):q_N(u)=q_N(v)\}.
\]
\begin{theorem}[Canonical difference filtration]\label{III:S:thm:filtration}
Each $U_N$ is transition-invariant, and
\[
 U_0=\SSpan\{s(w)-s(\Seps)\},\qquad
 U_N=\sum_{|q|=N}U_0A_q,\qquad
 U_{N+1}=\sum_aU_NA_a\subseteq U_N.
\]
The first plateau is permanent. If $\tau=\min\{N:U_N=U_{N+1}\}$, then $\tau\le\dim U_0-\dim U_\tau$. Also $U_N=0$ if and only if $f$ depends only on its $N$-suffix.
\end{theorem}
\begin{proof}
Equal windows remain equal after a common letter, proving invariance. Distinct words with a common full window have the form $xq,yq$, $|q|=N$, and their state difference is $(s(x)-s(y))A_q$. Shorter window fibers are singletons. This proves the closed form; group words by their last letter for the recursion. Nesting follows directly from equality of longer suffixes implying equality of shorter ones. A plateau persists under the same recursion, and each preceding proper inclusion loses at least one dimension. If $U_N=0$, outputs agree on fibers. Conversely if current outputs factor through the window, then outputs also agree after every continuation; column-span minimality makes the state difference zero.
\end{proof}

The initial space is the invariant closure of the finitely many rows $\alpha^\top(A_a-I)$. Indeed
\[
 s(a_1\cdots a_t)-\alpha^\top
 =\sum_{i=1}^t\alpha^\top(A_{a_i}-I)A_{a_{i+1}}\cdots A_{a_t}.
\]
Thus the entire filtration is computable by exact basis closure and row reduction, without enumerating all histories.

\begin{definition}[Charged affine interface]\label{III:S:def:affine-interface}
An order-$N$ implementation has a recurrent row $z(w)\in\SK^{1\times r}$ and
\[
 z(wa)=z(w)B_a+g(q_N(w),a),\qquad
 f(w)=z(w)c+h(q_N(w)).
\]
The matrices $B_a$ and recurrent readout $c$ are fixed; the forcing and direct-readout tables are unrestricted. The cost charges recurrent dimension and the actual unital algebra generated by $B_a$, while treating the window and tables as externally supplied. It does not count table bytes, arithmetic, or numerical precision.
\end{definition}

\begin{theorem}[Exact residual minimum and marked quotient]\label{III:S:thm:window}
Let $C_N$ be the image of the transition algebra on $U_N$, with $C_N=0$ when $U_N=0$. Within the stated interface the minimum recurrent dimension is $\dim U_N$, and the minimum homogeneous Loewy length is $\SLL(C_N)$. Every exact competitor has a letter-marked generated-algebra surjection onto $C_N$. One construction attains all these lower bounds and every graded radical dimension simultaneously.
\end{theorem}
\begin{proof}
Choose a row-basis matrix $E_N$ for $U_N$, and define $C_{N,a}$ by $E_NA_a=C_{N,a}E_N$. The coordinate relation
\[
 z(w)E_N=s(w)-s(q_N(w))
\]
is well-defined. For $q'=\Ssuff_N(qa)$, the forcing is defined uniquely by
$g(q,a)E_N=s(q)A_a-s(q')$, which belongs to $U_N$. With $c=E_N\beta$ and $h(q)=s(q)\beta$, this gives an exact attaining implementation.

For any competitor, let $V_N$ span its state differences within window fibers. Future windows match, so forcing and direct readout cancel:
\[
 f(ut)-f(vt)=(z(u)-z(v))B_tc.
\]
Send each competitor difference to $s(u)-s(v)$. A linear relation among competitor differences gives a target row annihilating every $A_t\beta$, hence zero by column-span minimality. The map $V_N\twoheadrightarrow U_N$ is therefore well-defined, onto, and intertwines all letters. This proves the dimension bound and a representation-kernel inclusion, hence the generated-algebra quotient. Apply Theorem~\ref{I:S:thm:marked} for the radical-layer bounds.
\end{proof}


\begin{theorem}[Saturation, one-symbol bound, and invertible padding]\label{III:S:thm:windowbounds}
Writing $\lambda_N=\SLL(C_N)$, one has
\[
 0\le\lambda_N-\lambda_{N+1}\le1,\qquad
 U_NJ(\ScA_f)^j\subseteq U_{N+j}.
\]
For nonzero $f$, $0\le\SLL(\ScA_f)-\lambda_0\le1$. If one transition is invertible, $U_N=U_0$ for every $N$. For a faithful finite-dimensional right $A$-module $M$, $L=\SLL(A)$, and $0\le r\le L$, the image algebra on $MJ(A)^r$ has Loewy length exactly $L-r$.
\end{theorem}
\begin{proof}
Restriction gives epimorphisms $C_N\twoheadrightarrow C_{N+1}$. The positive-word ideal $A_+=\SSpan\{A_w:|w|\ge1\}$ has quotient zero or $\SK$, so $J(A)\subseteq A_+$. Invariance and the recursion give $U_NA_+=U_{N+1}$, hence $U_NJ(A)\subseteq U_{N+1}$. Iteration proves the inclusion, and if $U_{N+1}J(A)^k=0$, then $U_NJ(A)^{k+1}=0$. This proves the one-symbol bound, including the zero endpoint. The quotient $M/U_0$ is at most one-dimensional and every letter acts as the identity; it is semisimple. Thus $MJ(A)\subseteq U_0$, giving the centering bound by faithfulness. If $A_a$ is invertible, its restriction to invariant $U_0$ is onto; hence $U_1\supseteq U_0A_a=U_0$ and the plateau persists. Finally, on $MJ^r$ the $k$th radical power of the image vanishes exactly when $MJ^{r+k}=0$, or equivalently $J^{r+k}=0$ by faithfulness. The least such $k$ is $L-r$, with the zero endpoint treated separately.
\end{proof}


\begin{example}[The neutral-letter counterexample]
For $f(w)=\mathbf1\{\#a(w)=1\}$ over $\{a,b\}$, take $A_a=\begin{psmallmatrix}0&1\\0&0\end{psmallmatrix}$, $A_b=I$, initial row $(1,0)$ and final column $(0,1)^\top$. The rank and Loewy length are two, and $U_0=\SR^2$. The invertible neutral letter forces $U_N=\SR^2$ at every order. Histories $b^N$ and $ab^N$ already have the same raw window and different current outputs. Finite algebraic nilpotence in one letter does not imply a saving from a raw token window when another letter can pad indefinitely.
\end{example}

\begin{example}[Sharp decrease and the limits of radical tails]
The unary length-$m$ impulse has $\dim U_N=\lambda_N=\max(m+1-N,0)$, since $U_0$ is the whole nilpotent shift space and $U_N=U_0A^N$. This attains the one-symbol bound. Adding $(-1)^{|w|}$ preserves a semisimple coordinate in every $U_N$: the dimension becomes $\max(m+1-N,0)+1$ and the Loewy cost is $\max(m+1-N,1)$. No positive radical tail contains that semisimple coordinate. Thus the exact radical-tail restriction theorem does not assert that such tails optimize every window problem.
\end{example}

\subsection{Window-conditioned homogeneous dynamics}
If the homogeneous matrices and recurrent readout may also depend on the window, define fibers
\[
 F_N(q)=\SSpan\{s(u)-s(v):q_N(u)=q_N(v)=q\}.
\]
They are zero for $|q|<N$, and equal $U_0A_q$ for $|q|=N$. Letters map $F_N(q)$ into $F_N(\Ssuff_N(qa))$. Let $\mathcal P_N$ be the image of the resulting path algebra on $\bigoplus_q F_N(q)$, including the vertex idempotents.

\begin{definition}[Window-indexed affine interface]
\label{III:S:def:fiber-interface}
Let $\mathcal Q_N=\{q_N(w):w\in\Sigma^*\}$, with successor
$q'=q_N(qa)$. Choose a dimension $r_q\ge0$ for each vertex $q$.
The recurrent row after history $w$ belongs to
$z(w)\in\mathbb K^{1\times r_{q_N(w)}}$, and the update/readout are
\[
 z(wa)=z(w)B_{q,a}+g_{q,a},\qquad
 f(w)=z(w)c_q+h_q,\qquad q=q_N(w),
\]
where $B_{q,a}\in\mathbb K^{r_q\times r_{q'}}$,
$g_{q,a}\in\mathbb K^{1\times r_{q'}}$ and
$c_q\in\mathbb K^{r_q\times1}$. The active window is external.
We charge peak recurrent dimension $\max_q r_q$ and, separately,
the Loewy length of the path-image algebra on the direct sum of
within-window reachable difference spaces, with vertex idempotents
and letter arrows marked. Window and forcing-table storage is not
included. This extends Definition~\ref{III:S:def:affine-interface} by
allowing the homogeneous update and readout to depend on the window.
\end{definition}

\begin{theorem}[Fibered minimum]\label{III:S:thm:fiber}
Within Definition~\ref{III:S:def:fiber-interface}, the minimum peak recurrent dimension is $\max_q\dim F_N(q)$. The minimum Loewy length of the actual marked path-image algebra is $\SLL(\mathcal P_N)$. These are different objectives from Theorem~\ref{III:S:thm:window}.
\end{theorem}
\begin{proof}
Within each fiber, the difference map used in Theorem~\ref{III:S:thm:window} is well-defined and onto, since future window paths and their forcing agree. The direct sum of these maps intertwines every arrow and every vertex idempotent, giving a path-algebra quotient and its radical lower bounds. Choosing bases for the canonical fibers attains the bounds, with translations supplied by the reference rows $s(q)$. Only one vertex is active in an execution, so the physical dimension charged here is the maximum, although the algebra acts on a direct sum. If desired, each fiber can be embedded into that common peak dimension; unused coordinates are set to zero.
\end{proof}

These different algebra costs need not decrease together. For the neutral-letter example at $N=1$, order the fiber coordinates as one coordinate at $a$ and two at $b$. The path radical is spanned by $E_{13},E_{21},E_{23}$; its square contains $E_{23}$ and its cube is zero. Its Loewy length is three, although the fixed-letter cost is two. This is not an extra saving or an extra physical register forced by the fixed-letter theorem: the marked algebra being charged has changed.

\section{Finite-energy forgetting on the difference module}\label{III:S:sec:metric}
For a real finite realization define the unnormalized future-energy diameter
\[
 \omega_N(f)=\sup_{q_N(u)=q_N(v)}
       \left(\sum_{t\in\Swords}|f(ut)-f(vt)|^2\right)^{1/2},
\]
allowing $+\infty$. It is nonincreasing in $N$ and vanishes exactly for suffix-definite behavior. If $H_f,H_g$ are bounded, then
\begin{equation}\label{III:S:eq:omegalipschitz}
 |\omega_N(f)-\omega_N(g)|\le\sqrt2\Snorm{H_f-H_g}.
\end{equation}
Indeed each row difference is tested by $e_u-e_v$, of norm $\sqrt2$ when $u\ne v$. The constant is sharp: a unary impulse of length $m\ge N+1$ has operator norm one and two such residual rows at distance $\sqrt2$.

Choose a basis matrix $E$ of $U_0$, and define $EA_a=C_aE$, $b=E\beta$, $g_aE=\alpha^\top(A_a-I)$. Center states by $s(w)=\alpha^\top+z(w)E$. Then
\[
 z(wa)=z(w)C_a+g_a,\qquad z(xq)=z(x)C_q+z(q).
\]
\begin{theorem}[Correct analytic domain for forgetting]\label{III:S:thm:diffstable}
For a minimal finite realization with $U_0\ne0$, these are equivalent:
\[
 \omega_0<\infty;\quad
 G_U=\sum_tC_tbb^\top C_t^\top\text{ finite};\quad
 \rho(\mathcal T)<1;\quad
 \exists R\succ0:R-\mathcal T(R)\succ0,
\]
where $\mathcal T(X)=\sum_aC_aXC_a^\top$. Then $G_U\succ0$ and
\[
 \omega_N^2=\max_{|q|=N}\sup_{x,y}
 (z(x)-z(y))C_qG_UC_q^\top(z(x)-z(y))^\top.
\]
No boundedness of the full $H_f$ is assumed. For $U_0=0$ all diameters vanish.
\end{theorem}
\begin{proof}
Choose a basis from the realized state differences spanning $U_0$. Finite $\omega_0$ makes each of its future sequences square summable, and finite-dimensional linearity gives a finite Gramian. Minimal future observability on $U_0$ makes it positive definite. Its tail identity $\mathcal T^j(G_U)\to0$ and the order-unit argument of Theorem~\ref{I:S:thm:stability} prove strict positive-map stability. Conversely strict stability yields $R=I+\mathcal T(R)$. Every $C_a$ is a strict contraction in the row norm $\Snorm x_R^2=xRx^\top$, with a common factor below one. The affine forcing belongs to a finite set, so all centered states are uniformly bounded by a geometric series. The finite Gramian then bounds their future diameter. Finally a common full window $q$ propagates a difference by $C_q$, proving the formula.
\end{proof}


\begin{theorem}[Finite computation with a certified remainder]\label{III:S:thm:finiteomega}
Choose $R\succ0$, $0\le\eta<1$, $c>0$ such that $C_aRC_a^\top\preceq\eta^2R$ and $G_U\preceq cR$. Put $M=\max_a\Snorm{g_a}_R/(1-\eta)$ and
\[
 D_{N,L}=\max_{\substack{|q|=N\\|x|,|y|\le L}}
             \Snorm{(z(x)-z(y))C_q}_{G_U}.
\]
Then
\[
 \omega_N\le2\sqrt c M\eta^N,\qquad
 D_{N,L}\le\omega_N\le D_{N,L}+2\sqrt c M\eta^{N+L}.
\]
With rational data this computes $\omega_N$ to arbitrary positive accuracy by finite exact arithmetic and certified bounds on square roots.
\end{theorem}
\begin{proof}
Induction gives $\Snorm{z(w)}_R\le M$. Replacing a long history by its last $L$ letters changes its centered state by at most $M\eta^L$. Replace both histories and propagate the common $N$-word, then dominate the future norm by $\sqrt c\Snorm\cdot_R$. This proves the remainder. Every finite pair is an admissible pair for the infinite supremum, proving the lower bound. Exact Stein systems give rational $R,G_U$; rational matrix inequalities certify upper constants, and $D_{N,L}^2$ is a finite maximum of rational numbers. Use $\eta^0=1$, including $\eta=0$.
\end{proof}


\begin{theorem}[Exact root rate and finite extinction]\label{III:S:thm:jsr}
Under Theorem~\ref{III:S:thm:diffstable}, let $r=\widehat\rho(\{C_a\})$ be the joint spectral radius. Then $0\le r<1$ and $\lim_N\omega_N^{1/N}=r$. If $0<r<1$ and $N_\epsilon=\min\{N:\omega_N\le\epsilon\}$, then
\[
 \lim_{\epsilon\downarrow0}\frac{N_\epsilon}{\log(1/\epsilon)}
  =\frac1{\log(1/r)}.
\]
For any finite-dimensional real or complex difference family, $r=0$ is equivalent to simultaneous vanishing of all sufficiently long products, to $U_N=0$ for some $N$, and to finite suffix-definiteness. Extinction occurs by $\dim U_0$ when nonzero.
\end{theorem}
\begin{proof}
Put $b_N=\max_{|q|=N}\Snorm{C_q}_{G_U}$. The defining state differences span, so choose a row basis $d_i$ among them. There is a fixed $\kappa$ such that $\Snorm T_{G_U}\le\kappa\max_i\Snorm{d_iT}_{G_U}$ for all $T$, by coordinate norm equivalence. Thus
$\kappa^{-1}b_N\le\omega_N\le\omega_0b_N$. Submultiplicativity gives $b_N^{1/N}\to r$, proving the rate. The common contraction gives $r<1$. Sandwich $\omega_N$ eventually between $r_-^N$ and $r_+^N$ for $r_-<r<r_+$ and use monotonicity to obtain the inverse logarithmic limit.

For extinction, $r=0$ implies every nonempty word product is nilpotent by repeating that word. In the generated algebra the positive-word ideal $A_+$ consists of nilpotent elements: traces of every power of $X\in A_+$ vanish because they are sums of traces of positive word products. Newton identities and Cayley--Hamilton imply $X^d=0$ in dimension $d$. For any $Y\in A$, $YX\in A_+$, so $I-YX$ is invertible. The radical criterion gives $A_+\subseteq J(A)$, hence $A_+$ is a nilpotent ideal. All sufficiently long products vanish. The descending chain $W_{j+1}=\sum_aW_jC_a$, $W_0=\SK^{1\times d}$, cannot have a nonzero plateau before zero because that would contradict this nilpotence. There are at most $d$ strict drops. The equivalence with $U_N=0$ and suffix-definiteness follows from Theorem~\ref{III:S:thm:filtration}. The reverse implication to zero JSR is immediate.
\end{proof}

This rate comparison is a specialization of finite-dimensional realization and switched-system stability, not a new general theorem about the JSR. For any $r<\eta<1$ there is a norm $\Snorm x_\eta=\sup_{j,w:|w|=j}\eta^{-j}\Snorm{xC_w}_{G_U}$ with contraction factor $\eta$. A finite version suffices once a block length $m$ satisfies $C_wG_UC_w^\top\preceq\eta^{2m}G_U$ for every $|w|=m$: take the maximum only over lengths $0,\ldots,m-1$. Multiplying once shifts each term to the next length, and the last length is controlled by the block inequality. This need not be a quadratic norm.

\begin{example}[Persistent drive; exact memory versus metric forgetting]
Let $z(wa)=\rho z(w)+1$, $z(wb)=\rho z(w)-1$, $z(\Seps)=0$, with $0<\rho<1/\sqrt2$ and $f=z$. The supremum distance between states is $2/(1-\rho)$, and $G_U=1/(1-2\rho^2)$. Hence
\[
 \omega_N=\frac{2\rho^N}{(1-\rho)\sqrt{1-2\rho^2}}.
\]
The full Hankel is unbounded since $f(a^n)\to1/(1-\rho)\ne0$. The homogeneous difference dimension remains one despite metric forgetting. In contrast, $C_a=C_b=3/4$ has JSR $3/4<1$ but all-word positive-map radius $9/8>1$. The JSR stability domain cannot replace the finite-energy hypothesis. A Jordan block at eigenvalue $1/2$ has powers with a factor $N2^{-N}$; a root rate of $1/2$ does not imply a bound $K2^{-N}$.
\end{example}

\section{Conditional futures: normalization changes the question}\label{III:S:sec:conditional}
\subsection{Completed-string probability semantics}
Assume $\alpha,\beta,A_a$ are nonnegative, $T=\sum_aA_a$ has $\rho(T)<1$, and
\[
 v=(I-T)^{-1}\beta\text{ is entrywise positive},\qquad \alpha^\top v=1.
\]
Then $f(w)=\alpha^\top A_w\beta$ is a distribution on \emph{completed finite strings}. For an admissible prefix $u$, define
\[
 F(u)=s(u)v=\sum_t f(ut)>0,\quad
 \bar s(u)=s(u)/F(u),\quad p(t\mid u)=\bar s(u)A_t\beta.
\]
The vectors $\bar s(u)$ lie in the compact simplex $\mathcal S_v=\{r\ge0:rv=1\}$. The future Gramian $Q=\sum_tA_t\beta\beta^\top A_t^\top$ is finite, since each coordinate future sequence has finite $\ell^1$ mass $v_i$. It is positive definite precisely when the future columns span. If $V=\Sdiag(v)$, the gauge $r\mapsto rV$, $A_a\mapsto V^{-1}A_aV$, $\beta\mapsto V^{-1}\beta$ gives $\widetilde\beta+\sum_a\widetilde A_a\mathbf1=\mathbf1$.

Define
\[
 \widehat\omega_N^2=
 \sup_{\substack{q_N(u)=q_N(w)\\F(u),F(w)>0}}
       (\bar s(u)-\bar s(w))Q(\bar s(u)-\bar s(w))^\top.
\]
It is a distance between probability laws, so $0\le\widehat\omega_N\le\sqrt2$, it is nonincreasing, and it vanishes exactly when conditional complete futures agree within each admissible window fiber. In an observable realization this is equality of normalized states; it is not $U_N=0$ for unnormalized states. Cylinder probabilities of nonterminating processes require a separately declared horizon or weighting; their unrestricted sum-of-squares metric is not substituted here.

\begin{theorem}[Exact projective difference identity]\label{III:S:thm:projective}
Let $\Psi_q(r)=rA_q/(rA_qv)$ on its positive-mass domain and $B(y)=I-vy$ for $yv=1$. Then $B(y)^2=B(y)$, and
\[
 \Psi_q(r)-\Psi_q(r')
 =\frac{(r-r')A_qB(\Psi_q(r'))}{rA_qv},\qquad
 D\Psi_q(r)[h]=\frac{hA_qB(\Psi_q(r))}{rA_qv}\quad(hv=0).
\]
Equivalently, the numerator can use $(s(x)-s(y))A_q$ and the denominator $F(xq)$ for histories $xq,yq$.
\end{theorem}
\begin{proof}
For $y=\Psi_q(r')$, $r'A_qB(y)=0$ and $rA_qB(y)=rA_q-(rA_qv)y$. Divide by $rA_qv$. The idempotence follows from $yv=1$. Differentiate the exact identity or the normalized quotient to obtain the derivative.
\end{proof}


\begin{theorem}[Finite row-simplex upper certificate]\label{III:S:thm:rowsimplex}
For $|q|=N$, let $b_i(q)=e_i^\top A_qv$ and $r_i(q)=e_i^\top A_q/b_i(q)$ for live rows with $b_i(q)>0$. Then
\[
 \widehat\omega_N^2\le
 \max_{|q|=N}\max_{i,j:\ b_i(q)b_j(q)>0}
           (r_i(q)-r_j(q))Q(r_i(q)-r_j(q))^\top.
\]
The right side is exactly the squared diameter when arbitrary inputs from $\mathcal S_v$ are allowed. It need not equal the reachable-history diameter. All its squared terms are rational for rational model data.
\end{theorem}
\begin{proof}
For any admissible $r$,
\[
 \Psi_q(r)=\sum_{i:b_i(q)>0}\frac{r_i b_i(q)}{rA_qv}\,r_i(q).
\]
This is a convex combination. A zero row mass implies a zero row of $A_q$ because $v>0$. The diameter of a convex hull in the $Q$ seminorm equals the largest vertex distance, by expanding the difference of two mixtures into a convex combination of pairwise differences. Each live vertex is reached from $e_i^\top/v_i\in\mathcal S_v$. If the requisite vertices lie in the closure of the normalized reachable set, continuity at live vertices also yields equality for reachable histories. The unique rational Stein system supplies rational $Q$.
\end{proof}


\begin{proposition}[A qualified mass-floor bound]
Let $D=\operatorname{diam}_2(\mathcal R)$ for the normalized reachable set, $\Pi=\sup_{y\in\mathcal S_v}\Snorm{I-vy}_2$, $a_N=\max_{|q|=N}\Snorm{A_q}_2$, and $m_N=\inf_{r\in\mathcal R,\ |q|=N,\ rA_qv>0}rA_qv$. If $m_N>0$, then
\[
 \widehat\omega_N\le\min\{\sqrt2,\sqrt{\lambda_{\max}(Q)}\Pi D a_N/m_N\}.
\]
\end{proposition}
\begin{proof}
Apply Theorem~\ref{III:S:thm:projective} and norm submultiplicativity. Compactness gives $\Pi<\infty$. A zero mass infimum makes this particular bound vacuous, not false or divergent.
\end{proof}


\subsection{The precise unary conditional exponent}
For a fixed letter $j$, let $\widehat\omega_{j,N}$ be the supremum over histories $xj^N,yj^N$, rather than over every length-$N$ window.
\begin{theorem}[Perron rate with visible modes and uniform positive overlap]\label{III:S:thm:perron}
Assume $A=A_j$ is primitive. Let $\lambda>0$ be its Perron root and choose positive $h,\ell$ with $Ah=\lambda h$, $\ell A=\lambda\ell$, $\ell h=1$. Put $P=h\ell$, $R=A/\lambda-P$, and $\theta=\rho(R)<1$. For $0<\theta<1$, let $b$ be the largest Jordan-block size on the spectral circle $|z|=\theta$. Then
\[
 \widehat\omega_{j,N}\le K(1+N)^{b-1}\theta^N.
\]
If the realization has both full row and full column spans, then
$\lim_N\widehat\omega_{j,N}^{1/N}=\theta$. If $\theta=0$, sufficiently long $j$-runs have identical normalized futures. Without the span assumptions the upper bound remains valid, but equality of the observable rate need not hold.
\end{theorem}
\begin{proof}
For $N\ge1$, $(A/\lambda)^N=P+R^N$. Put $t_*=\ell/(\ell v)$. Direct normalization yields
\[
 \Psi_{j^N}(r)-t_*=
 \frac{rR^N(I-vt_*)}{(rh)(\ell v)+rR^Nv}.
\]
On the compact simplex, $rh\ge\min_i h_i/v_i>0$. Thus denominators are uniformly positive for all large $N$, and the finitely many smaller $N$ cause no problem because $A^Nv>0$. The normalized-state diameter is bounded by a constant times $\Snorm{R^N}$. The Jordan power formula gives the asserted polynomial-exponential upper estimate.

For the lower estimate choose independent reachable normalized rows $r_1,\ldots,r_n$, possible under full row span. Set $w_i=r_i/(r_ih)$ and $d_i=w_i-w_1$ for $i\ge2$. These form a basis of $\{x:xh=0\}$; together with $\ell$ they span all rows. For $u_{i,N}=\Psi_{j^N}(r_i)$ one has
\[
 d_iR^N=\frac{u_{i,N}}{u_{i,N}h}-\frac{u_{1,N}}{u_{1,N}h},\qquad\ell R^N=0.
\]
The map $u\mapsto u/(uh)$ is uniformly Lipschitz on the simplex. Full future span gives $Q\succ0$ and norm equivalence, so $\Snorm{R^N}\le K'\widehat\omega_{j,N}$. Taking roots and using the spectral-radius formula proves equality. At $\theta=0$, $R$ is nilpotent and the displayed normalized error eventually vanishes.
\end{proof}

This is a specialization of Perron--Frobenius and projective-filtering theory~\cite{cohen,atar}. A simple dominant eigenvalue alone is insufficient, and a Jordan factor cannot in general be removed. Subsection~\ref{III:S:app:conditional} gives exact rational counterexamples to both omissions.

\begin{example}[Exact separation of mass rate and conditional rate]\label{III:S:ex:rates}
Take $\alpha^\top=(1/2,1/2)$,
\[
 A_j=\begin{pmatrix}6/25&4/25\\4/25&6/25\end{pmatrix},\quad
 g=\binom{1/5}{1/10},\quad A_a=ge_1^\top,\quad A_b=ge_2^\top,
 \quad\beta=\binom{1/5}{2/5}.
\]
Then $v=\mathbf1$, the realization is minimal, and the reset letters reach both simplex vertices. $A_j=(2/5)P+(2/25)(I-P)$ for $P=\mathbf1\mathbf1^\top/2$. Exact Stein calculation gives
\[
 (1,-1)Q(1,-1)^\top=700/16369,\qquad
 \widehat\omega_{j,N}^2=(700/16369)25^{-N}.
\]
Thus the conditional root rate is $1/5$. The unnormalized pair $aj^N,bbj^N$ has a nonzero Perron contrast because $F(a)=3/20$, $F(bb)=3/200$, so its future-energy root rate is $2/5$. These are exact rates on an explicitly specified fixture, not inferred rates from an unavailable numerical trace.
\end{example}

\begin{theorem}[Uniform positive-matrix certificate]
If every $A_a$ is strictly positive, define
\[
 \Delta_a=\log\max_{i,j,k,l}\frac{(A_a)_{ik}(A_a)_{jl}}{(A_a)_{il}(A_a)_{jk}},\quad
 \Delta_* =\max_a\Delta_a,\quad\tau_* =\max_a\tanh(\Delta_a/4)<1.
\]
For $N\ge1$,
\[
 \widehat\omega_N\le
 \min\{\sqrt2,2\tanh(\Delta_*\tau_*^{N-1}/4)\}
 \le(\Delta_*/2)\tau_*^{N-1}.
\]
When $\Delta_*=0$, interpret the bound as zero after the first letter.
\end{theorem}
\begin{proof}
The classical Birkhoff theorem gives projective image diameter $\Delta_a$ and contraction factor $\tanh(\Delta_a/4)$~\cite{cohen}. The first positive update sends boundary states into the cone interior, and each following update contracts the Hilbert distance. Thus any $N$-word has diameter at most $\Delta_*\tau_*^{N-1}$. For probability rows the sharp comparison is $\Snorm{\mu-\nu}_1\le2\tanh(H(\mu,\nu)/4)$, with $\ell^1$ rather than half-$\ell^1$ convention~\cite{cohen}. The gauge by $V$ preserves projective distance. Complete-future laws are mixtures of state future laws, so their $\ell^2$ distance is at most the mixture-weight $\ell^1$ distance. Combine these facts and $\tanh x\le x$.
\end{proof}


\begin{theorem}[Mass-qualified transfer; failure without a margin]\label{III:S:thm:conditionaltransfer}
Let $f,g$ be completed-string distributions and use one fixed history set $\mathcal H$ admissible in both, with $F_f(u),F_g(u)\ge\mu>0$. If $\Snorm{H_f-H_g}\le e$ and $\sup_{u\in\mathcal H}|F_f(u)-F_g(u)|\le e_F$, then
\[
 |\widehat\omega_N^{\mathcal H}(f)-\widehat\omega_N^{\mathcal H}(g)|
 \le2(e+e_F)/\mu.
\]
The hypothesis $\Snorm{f-g}_1\le\epsilon$ suffices for the alternative bound $4\epsilon/\mu$. No unqualified uniform continuity in Hankel norm holds for conditional diameters.
\end{theorem}
\begin{proof}
For a history row $r_u^f(t)=f(ut)$,
\[
 \Snorm{r_u^f/F_f-r_u^g/F_g}_2
 \le\frac{\Snorm{r_u^f-r_u^g}_2+|F_f-F_g|}{F_f},
\]
because $\Snorm{r_u^g}_2\le F_g$. A Hankel row error has norm at most $e$. Use the bound for both histories and the reverse triangle inequality, then take the same supremum. An $\ell^1$ error bounds both subtree mass and row errors by $\epsilon$.

For failure without a margin, assign $f_\epsilon(b)=1-\epsilon$, $f_\epsilon(ac)=f_\epsilon(bad)=\epsilon/2$, and $g_\epsilon(b)=1-\epsilon$, $g_\epsilon(ac)=g_\epsilon(bac)=\epsilon/2$, with other masses zero. In $f_\epsilon$ the histories $a,ba$ have deterministic distinct futures, so $\widehat\omega_1(f_\epsilon)=\sqrt2$. In $g_\epsilon$ all admissible equal-last-letter histories have identical futures, so its diameter is zero. Yet $\Snorm{f_\epsilon-g_\epsilon}_1=\epsilon$ and Proposition~\ref{I:S:prop:normbounds} gives $\Snorm{H_{f_\epsilon}-H_{g_\epsilon}}\le\epsilon\to0$.
\end{proof}

Hankel norm alone does not bound $e_F$: summing an infinite residual against the all-ones vector is not a bounded $\ell^2$ functional. Nor does an $\ell^2$ bound on completed futures imply a dimension-free total-variation or KL bound. In a finite set of $M$ outcomes, $\Snorm{p-q}_1\le\sqrt M\Snorm{p-q}_2$; using this inequality on an infinite completion space would be invalid.

\section{Behavioral equivalence and autonomous certification}
\label{III:M:sec:behavioral-contract}
This section supplies elementary bridges needed to read the later results
as a theory of agent memory. They are classical arguments, not asserted
novelty. Fix an event alphabet $\Sigma$, a query set $\mathcal Q$ and a
common answer space. Assume every history and continuation used below is
admissible; for a restricted language, specify its state and common
admissibility conditions as additional persistent information.

\paragraph{From agent records to a predictive contract.}
For an agent-memory application, an implementation must specify how
observations, feedback, and corrections become letters, which queries
are admissible, and what answer law or scalar observable is to be
preserved. Source records and provenance can support validation of this
encoding. The external-memory organization discussed by
Martin~\citep{martin2026memory} motivates retaining enough evidence to
revise a mistaken rule. The mathematical replay statements below assume
that the supplied events and mode information describe the intended
updates; they do not by themselves establish that a natural-language
extractor identified the right rule or correction.

FrontierMemory's task, user, and environment cases suggest complementary
tests of this interface~\citep{letta2026frontiermemoryv1}. A proposed
DRAGA evaluation should freeze the event and query interfaces, separate
memory construction from held-out tasks, and compare task performance
with and without the retained memory. It should also measure update
correctness, correction of stale records, complete stored bytes, and
online and offline computation. This accounting matters because the
benchmark's reported task costs exclude its offline memory-generation
cost. These are evaluation requirements for a future implementation;
this paper reports no FrontierMemory score. The pseudometric below
instead states a worst-case preservation target for a specified model,
including the effect of future continuations.

\begin{definition}[Future answer pseudometric]\label{III:M:def:agent-future}
For reference answer laws $P_h(\cdot\mid q)$ define
\[
 d_{\rm fut}^{\rm ans}(h,h')
 =\sup_{w\in\Sigma^*,\,q\in\mathcal Q}
 \MTV\!\left(P_{hw}(\cdot\mid q),P_{h'w}(\cdot\mid q)\right).
\]
For scalar observables replace TV by absolute difference. The scalar
metric in the later memory theorems is this latter convention; a scalar
readout theorem is not automatically a theorem for an answer distribution.
\end{definition}

\begin{proposition}[Exact quotient and behavioral packing]
\label{III:M:prop:agent-quotient}
Zero future distance defines a right congruence:
$h\sim h'$ implies $ha\sim h'a$ for every letter. If it has finitely
many classes, its quotient is an exact autonomous predictive machine.
Every exact deterministic machine has at least that many reachable states.
More generally, if $M$ histories have pairwise future distances greater
than $2\epsilon$, every deterministic machine with uniform answer error at
most $\epsilon$ has at least $M$ reachable states and needs
$\lceil\log_2M\rceil$ persistent bits.
\end{proposition}
\begin{proof}
The supremum defining $d_{\rm fut}^{\rm ans}(ha,h'a)$ ranges only over
continuations beginning with $a$, so it is at most
$d_{\rm fut}^{\rm ans}(h,h')$. This proves the congruence and makes
the quotient transition $[h]\mapsto[ha]$ well defined. Empty continuation
shows that the answer law is constant on each class. Conversely, equal
machine states remain equal after every common continuation, so an exact
machine cannot merge distinct classes. For approximate machines,
if two selected histories shared a state, their reference answers after
any common continuation would be at distance at most $2\epsilon$ by the
triangle inequality. This contradicts their separating future.
\end{proof}

The quotient result is an existence/minimality statement. An infinite
supremum, exact class equality or the quotient transitions need not be
computable from a finite data sample.

\begin{proposition}[A sufficient finite invariant-cover certificate]
\label{III:M:prop:agent-cover}
Let reference states evolve by $x^+=F_a(x)$. Suppose finitely many
nonempty sets $C_i$, a start index $i_0$, transitions $T(i,a)$ and
answer laws $\widehat P_i(\cdot\mid q)$ satisfy
\[
 x_0\in C_{i_0},\qquad F_a(C_i)\subseteq C_{T(i,a)},\qquad
 \sup_{x\in C_i,q\in\mathcal Q}\MTV(P_x(\cdot\mid q),
                     \widehat P_i(\cdot\mid q))\le\epsilon.
\]
Then storing only the index $i$ and applying $T$ gives uniform answer
error at most $\epsilon$ for all histories.
\end{proposition}
\begin{proof}
Induction gives $x_t\in C_{i_t}$. The final displayed bound applies
to every query at every time. Transitions and answers use only the index
and current input, so the implementation is autonomous.
\end{proof}

The finite cover alone does not supply its transition inclusion or an
efficient description. In particular, the relation
$d_{\rm fut}\le\epsilon$ is generally not an equivalence relation:
three real scalar states $0,3\epsilon/4,3\epsilon/2$ already violate
transitivity. Pairwise similarity clustering is therefore not a proof
of an update-compatible quotient.

\begin{proposition}[Bounded-loss decision transfer]\label{III:M:prop:decision-transfer}
Let $P,Q$ be answer laws with $\MTV(P,Q)\le\epsilon$. For actions
$a$ and loss $\ell(a,y)\in[0,1]$, suppose minimizers $a_P,a_Q$
exist. Then
\[
 \mathbb E_P\ell(a_Q,Y)-\mathbb E_P\ell(a_P,Y)\le2\epsilon.
\]
If the $P$-optimal action is unique with loss gap greater than
$2\epsilon$ to every other action, it remains the unique $Q$-optimal action.
\end{proposition}
\begin{proof}
Expectations of any $[0,1]$-valued function differ by at most TV.
Insert and subtract both $Q$ expectations. Their difference is
nonpositive by optimality of $a_Q$, leaving at most $2\epsilon$.
For the gap assertion, the difference between any two action risks
can change by at most $2\epsilon$.
\end{proof}

This result concerns the declared reference loss and query. It does not
convert a completed-word TV bound into a prefix-uniform guarantee, nor
does it certify that two agents generating different actions encounter
the same future observations.

\subsection{Where the multiletter aspect enters}
The order of events matters: $F_aF_b$ need not equal $F_bF_a$.
A direction invisible to the current readout may become observable
after another letter. A valid reduction must control all declared future
readouts and remain compatible with every allowed update.
Hankel theory encodes this past--future interaction linearly; the invariant
cover and stopped-address constructions handle finite executable states.
Neither linear Hankel rank nor spectral magnitude alone determines how
many persistent factual distinctions survive.

\paragraph{Convention for word order.}
The linear sections use $A_w=A_{a_1}\cdots A_{a_t}$ and a row state
updated on the right. Sequential geometric sections use
$x_t=\tau_{a_t}\circ\cdots\circ\tau_{a_1}(x_0)$.
The unequal-ratio address is written newest first and explicitly reversed
for chronological replay. These conventions describe the same input
ordering in their respective state orientations; they must not be mixed.

\section{An exact optimum in persistent bits}\label{III:M:sec:memory}
The following exact optimum does not depend on
the real dimension of the deployed representation.

\begin{definition}[Autonomous deterministic memory]\label{III:M:def:autonomous-memory}
A machine is a finite Moore machine $(Q,q_0,T_0,T_1,y)$ with deterministic,
time-independent transitions and output depending only on its current state.
All persistent information, including a clock, branch flag, cache, or
coroutine state, belongs to $Q$. Fixed code and immutable parameters may be
arbitrary, but cannot depend on the past input. Define
\[
 B_f^*(\epsilon)=\min\left\{\lceil\log_2|Q|\rceil:
       \sup_{w\in\{0,1\}^*}|f(w)-y(q(w))|\le\epsilon\right\}.
\]
An external clock or access to the deleted input history changes the model.
\end{definition}

Fix an integer $a\ge1$, put $\rho=2^{-a}$, and let
\[
 A_c=\begin{pmatrix}\rho&\rho c/2\\0&\rho/2\end{pmatrix},\quad
 \alpha=(1,0),\quad\beta=(0,1)^\top,\quad f_a(w)=\alpha A_w\beta.
\]
For a word $w=c_1\cdots c_t$, define $k(w)=\sum_{j=1}^tc_j2^{j-1}$.
The row-state recurrence and its closed form are
\[
 (z,x)\longmapsto(\rho z,\rho(x+cz)/2),\qquad
 z(w)=\rho^t,\quad f_a(w)=\frac{k(w)}{2^{(a+1)t}}.
\]
The letter matrices do not commute. The Hankel minor on
$\{\varepsilon,1\}$ has determinant $-\rho^2/4$, so rank is exactly
two. For $a=1$ its exact reachability and observability Gramians are
\[
 P=\begin{pmatrix}4/21&2/21\\2/21&8/7\end{pmatrix},\qquad
 Q=\begin{pmatrix}2&1/3\\1/3&4/21\end{pmatrix}.
\]
Thus $\|H_{f_1}\|_{\rm HS}^2=\operatorname{tr}(PQ)=292/441$.
Since $f_a(w)=(2\rho)^{|w|}f_1(w)$ and $2\rho\le1$, every source
has $\|H_{f_a}\|\le\sqrt{292/441}<5/6<1$.

\subsection{A finite witness for every future}
\begin{lemma}[Two suffixes determine the future metric]\label{III:M:lem:future}
For two states with differences $\Delta z,\Delta x$ and $0<\rho\le1/2$,
their supremum difference over all common suffixes is
\[
 d_{\rm fut}=\max\{|\Delta x|,\tfrac\rho2|\Delta x+\Delta z|\}.
\]
\end{lemma}
\begin{proof}
For a suffix of length $\ell$, the difference equals
$(\rho/2)^\ell\Delta x+\rho^\ell r_v\Delta z$ with
$0\le r_v\le1-2^{-\ell}$. A linear function on this interval is
bounded in absolute value by its endpoints. The first endpoint is at
most $|\Delta x|$. For $\ell\ge2$, the other can be written
$c_\ell(\rho/2)(\Delta x+\Delta z)-d_\ell\Delta x$ with
\[
 c_\ell=2\rho^{\ell-1}(1-2^{-\ell}),\qquad
 d_\ell=\rho^\ell(1-2^{1-\ell}),\qquad
 c_\ell+d_\ell\le5\,2^{-\ell}-6\,4^{-\ell}\le7/8.
\]
Lengths zero and one give the two asserted terms and attain them.
\end{proof}


\begin{lemma}[Behavioral packing]
If $M$ reachable prefixes have pairwise future distances greater than
$2\epsilon$, every admissible machine of error at most $\epsilon$
has at least $M$ reachable states.
\end{lemma}
\begin{proof}
Two merged prefixes give the same machine output after each common
suffix. Their exact outputs could then differ by at most $2\epsilon$,
contradicting the separating suffix.
\end{proof}

The packing principle is classical. The explicit construction and matching
integer bit count below are the specific result.

\subsection{The exact lower and upper bounds}
\begin{theorem}[Exact persistent-bit optimum]\label{III:M:thm:bits}
For $a\ge1$, $m\ge a+2$, and $\epsilon=2^{-((a+1)m+2)}$,
\[
 \boxed{B_{f_a}^*(\epsilon)=m+2.}
\]
The assertion is exact for $B=\lceil\log_2|Q|\rceil$, not for the minimum number of states. The lower bound allows arbitrary immutable tables. The upper construction
uses a closed integer update and includes its clock in the state identifier.
\end{theorem}
\begin{proof}[Lower bound]
Put $L=(a+1)m+2$ and take the prefix set
$\mathcal W_m=\{w:|w|\le m\}\cup\{0^{m+1},1^{m+1}\}$,
of size $2^{m+1}+1$. At one depth $t\le m$, distinct outputs differ
by at least $2^{-(a+1)t}\ge4\epsilon$.
The added pair has output gap
$2\rho(2^{m+1}-1)\epsilon\ge15\epsilon$.
At different depths, $|\Delta z|\ge\rho^m(1-\rho)$.
Lemma~\ref{III:M:lem:future} implies
$|\Delta z|\le(1+2/\rho)d_{\rm fut}$, hence
\[
 \frac{d_{\rm fut}}\epsilon\ge
 \rho 2^{m+2}\frac{1-\rho}{2+\rho}
 \ge16\frac{1-\rho}{2+\rho}\ge\frac{16}5.
\]
Every pair is separated by at least $16\epsilon/5>2\epsilon$.
Packing gives $|Q|\ge2^{m+1}+1$, and therefore at least $m+2$ bits.
\end{proof}

\begin{proof}[Upper bound]
Let $T=\lceil(L-1)/a\rceil$ and
$n_t=\min\{t,L-1-at\}$ for $0\le t<T$.
Create states $(t,b)$ for $0\le b<2^{n_t}$, plus an absorbing state
$\perp$. The integer $b$ records the most recent $n_t$ bits, newest
first. Encode the pair by consecutive blocks in a single identifier;
the clock is not stored separately. The initial state is $(0,0)$.
If $t+1<T$, on symbol $c$ update by
\[
 b'=c2^{n_{t+1}-1}
       +\left\lfloor b/2^{n_t-n_{t+1}+1}\right\rfloor,
 \qquad(t,b)\mapsto(t+1,b').
\]
All shifts have nonnegative integer exponents. The final layer enters
$\perp$, which loops forever. Output
\[
 y(t,b)=\begin{cases}b/2^{(a+1)t},&t\le m,\\
                      (2b+1)2^{-L},&m<t<T,
        \end{cases}\qquad y(\perp)=\epsilon.
\]
The update retains precisely the stated suffix. Before $m$ all bits are
retained. Later the discarded normalized binary tail is in
$[0,2^{-n_t}]$, so the midpoint error is at most
$\rho^t2^{-n_t-1}=\epsilon$. At and after $T$, the source is in
$[0,\rho^T)\subset[0,2\epsilon]$, so the absorbing output is valid.
This proves the all-horizon contract.

The constructed state count is $N=1+\sum_{t=0}^{T-1}2^{n_t}$.
The sum through $m$ equals $2^{m+1}-1$. For $t=m+j$, $j\ge1$,
one has $n_t=m+1-aj$, so the remaining finite geometric tail gives
\[
 N<2^{m+1}+\frac{2^{m+1}}{2^a-1}\le2^{m+2}.
\]
A single $(m+2)$-bit identifier therefore suffices. When $a=1$ the
exact constructed count is $2^{m+2}-2$.
\end{proof}

By monotonicity between successive tolerances, for each fixed $a$,
\[
 B_{f_a}^*(\epsilon)=\frac1{a+1}\log_2(1/\epsilon)+O(1).
\]
The coefficient varies while Hankel rank stays two. Thus rank alone does
not specify the optimal leading bit cost.

\begin{corollary}[Robustness and rank-four sources]\label{III:M:cor:robust-bits}
Under Theorem~\ref{III:M:thm:bits}, if $\|g-f_a\|_\infty\le\eta\le\epsilon/4$,
then $B_g^*(\epsilon+\eta)=m+2$. In particular
$B_g^*(5\epsilon/4)=m+2$ for every such $g$.
An explicit normalized rank-four example is
\[
 g(w)=f_a(w)+\frac\eta2(3^{-|w|}+5^{-|w|}),\qquad\eta=\epsilon/4.
\]
\end{corollary}
\begin{proof}
The same machine gives error $\epsilon+\eta$. Future distances decrease
by at most $2\eta$, leaving at least $16\epsilon/5-2\eta>
2(\epsilon+\eta)$. At tolerance $5\epsilon/4$, the lower margin is
at least $27\epsilon/10>5\epsilon/2$.

For the rank-four example append the scalar blocks $1/3,1/5$ to $A_c$.
Reachability rows and observability columns indexed by
$\varepsilon,0,00,1$ have full rank: the first three yield Vandermonde
minors in the distinct decay rates, and the last supplies the missing
coordinate. The extra series $h$ satisfies
$\|H_h\|_{\rm HS}\le9/7+25/23$ by the triangle inequality, since
$\sum_{t\ge0}(t+1)(2/c^2)^t=(1-2/c^2)^{-2}$ for $c=3,5$.
Consequently $(\eta/2)\|H_h\|_{\rm HS}\le(191/161)\eta<(9/7)\eta$.
The factor $1/2$ is required in this norm estimate;
the conclusion is unchanged. Since $\eta\le1/1024$,
$\|H_g\|<5/6+(9/7)\eta<1$.
\end{proof}


\subsection{What this solves, and what it leaves}
The optimum is exact in \emph{bits}, not in state count: the latter lies
between $2^{m+1}+1$ and the constructed $N$. Jointly encoding time is
essential. Decoding the layer by scanning block sizes takes $O(T)$
iterations and avoids an exponential table. A signed 64-bit implementation of these formulas
imposes $L\le62$; larger parameters need multiword arithmetic. Temporary
space and output precision may exceed the persistent state size.

Most importantly, $f_a$
decays uniformly with input length and the machine eventually absorbs.
It is a sharp executable-memory example, not evidence of sustained
language-memory retention. The rank-four perturbation depends on the
tolerance through $\eta$; it does not establish the same law for every
fixed rank-four source.
\subsection{A constant-overhead existence law on regular attractors}
\label{III:M:sec:ahlfors-memory}
For the following theorem, extend Definition~\ref{III:M:def:autonomous-memory}
to any fixed finite alphabet, with one deterministic transition for each
letter. Immutable transition and output tables are allowed and charged
separately from the persistent state identifier.

\begin{theorem}[Persistent-bit existence law under Ahlfors regularity]
\label{III:M:thm:ahlfors-memory}
Let $\mathcal X\subset\mathbb R^k$ be compact with positive diameter,
$x_0\in\mathcal X$, and maps $\tau_a:\mathcal X\to\mathcal X$ of
Lipschitz constant at most $0<\lambda<1$. Suppose the reachable states
are dense in $\mathcal X$. For $f(w)=h(\tau_w(x_0))$, assume
\[
 c_{\rm obs}\|x-y\|_2\le d_{\rm fut}(x,y)
 :=\sup_v|h(\tau_v(x))-h(\tau_v(y))|
 \le C_{\rm obs}\|x-y\|_2,
\]
with $0<c_{\rm obs}\le C_{\rm obs}<\infty$.
Let a finite Borel measure $\nu$ on $\mathcal X$ have total mass
$M\in(0,\infty)$ and satisfy, for open Euclidean balls and
$0<r<\operatorname{diam}\mathcal X$,
\[
 c_{\rm AD}r^s\le\nu(B(x,r)\cap\mathcal X)
                  \le C_{\rm AD}r^s,
 \qquad s,c_{\rm AD},C_{\rm AD}>0.
\]
Then, for $0<\varepsilon<c_{\rm obs}\operatorname{diam}\mathcal X/4$,
\begin{equation}
 s\log_2(1/\varepsilon)+C_-
 \le B_f^*(\varepsilon)
 \le s\log_2(1/\varepsilon)+C_+,
 \label{III:M:eq:ahlfors-bit-law}
\end{equation}
where
\[
 C_-=\log_2\!\left[\frac{M}{C_{\rm AD}}
                         \left(\frac{c_{\rm obs}}4\right)^s\right],
 \qquad
 C_+=1+\log_2\!\left[\frac{M}{c_{\rm AD}}
                   \left(\frac{2C_{\rm obs}}{1-\lambda}\right)^s\right].
\]
This is an existence theorem for persistent state bits. It does not
assert a computable net, a succinct immutable program, or an efficient
online state update.
\end{theorem}
\begin{proof}
Put $r_-=4\varepsilon/c_{\rm obs}$. A maximal finite
$r_-$-separated set has open $r_-$-balls covering $\mathcal X$;
maximality and compactness ensure both statements. Upper regularity
therefore gives at least
$M/(C_{\rm AD}r_-^s)$ centers. Their pairwise future distances are
at least $4\varepsilon$. By reachable density and the upper
Lipschitz bound, choose reachable states arbitrarily near the finitely
many centers so their future distances remain strictly greater than
$2\varepsilon$. Behavioral packing forces at least as many distinct
machine states and proves the lower bound.

For the upper bound, set
$\eta=(1-\lambda)\varepsilon/C_{\rm obs}$.
Choose a maximal $\eta$-separated set
$\{\widehat x_1,\ldots,\widehat x_N\}\subset\mathcal X$.
It is an $\eta$-net, and its open balls of radius $\eta/2$ are
pairwise disjoint. Lower regularity gives
\[
 N\le\frac{M}{c_{\rm AD}}
                   \left(\frac2\eta\right)^s.
\]
The net is chosen by separation; a minimum-cardinality net would not
automatically have the disjointness property used here.
Initialize at a nearest center to $x_0$, send center $i$ under letter
$a$ to a nearest center to $\tau_a(\widehat x_i)$, and output
$h(\widehat x_i)$. Tie-breaking is fixed. The tracking error satisfies
\[
 e_0\le\eta,\qquad e_{t+1}\le\lambda e_t+\eta,
 \qquad e_t\le\frac\eta{1-\lambda}
                          =\frac\varepsilon{C_{\rm obs}}.
\]
Since the empty suffix is allowed in $d_{\rm fut}$, output error is
at most $C_{\rm obs}e_t\le\varepsilon$. Finally
$\lceil\log_2N\rceil\le\log_2N+1$ yields $C_+$.
\end{proof}


The invariant-attractor condition
$\mathcal X=\bigcup_a\tau_a(\mathcal X)$ implies the required density
when $x_0\in\mathcal X$: every point lies in an iterated image of
diameter at most $\lambda^n\operatorname{diam}\mathcal X$, and the
same image of $x_0$ is reachable. The theorem's content is the
combination of packing and an invariant rounding tube; no novelty is
claimed for these standard ingredients. If finite-precision outputs are
required, reserve half the error for output quantization and halve the
tracking tolerance. This changes only the constant overhead, while all
output constants and transition tables still need to be charged.


\subsection{An optimal table-free executable on homogeneous separated attractors}
\label{III:M:sec:address-memory}
\begin{theorem}[Homogeneous addresses with all storage and update costs charged]
\label{III:M:thm:address-memory}
Suppose $\mathcal X\subset\mathbb R^k$ is the attractor of $d\ge2$
similarities $\tau_a(x)=\lambda O_ax+t_a$ of a common ratio
$0<\lambda<1$, with pairwise separated images. Let the dimension
$k$ and the rational affine maps be fixed. Assume $\mathcal X$ has
positive diameter, $x_0$ is the fixed point of $\tau_0$, and the
future-observability inequalities of Theorem~\ref{III:M:thm:ahlfors-memory}
hold for an output $h$ admitting certified $b$-bit evaluation in
polynomial time and $O_H(b)$ scratch bits. Then there is a table-free
finite-state executable with error at most $2^{-p}$ and
\[
 B(p)=\frac{\log d}{\log(1/\lambda)}p+O_H(1)
       =B_f^*(2^{-p})+O_H(1)
\]
persistent state bits. Its symbol update costs $O_H(p)$ bit operations.
For fixed-dimensional rational affine maps, readout takes
$O_H(p\mathsf M(p))$ bit operations plus the cost of evaluating $h$,
with $O_H(p)$ scratch and immutable parameter bits. No exponential
transition table is stored.
\end{theorem}
\begin{proof}
Put $D=\operatorname{diam}\mathcal X$ and choose
\[
 m=\max\left\{1,
 \left\lceil\frac{p+\log_2(2C_{\rm obs}D)}{\log_2(1/\lambda)}
                                      \right\rceil\right\}.
\]
The persistent state is a word $u$ of length $m$, initialized to $0^m$;
on a letter append it and discard the oldest symbol. Its mathematical
representative is $x(u)=\tau_{u_m}\circ\cdots\circ\tau_{u_1}(x_0)$.
Before $m$ inputs, the initial zeros fix $x_0$, so this representative
equals the true state. After $m$ inputs both states share the same
last $m$ maps, giving distance at most $D\lambda^m$ and output
error at most $2^{-p-1}$.

Pack the word as a base-$d$ integer $I\in\{0,\ldots,d^m-1\}$.
The exact update is $I\leftarrow(dI+a)\bmod d^m$.
It uses $\lceil m\log_2d\rceil$ persistent bits, without a startup
counter; $m,d^m$ and the program are immutable. Fixed-$d$ multiplication,
comparison and subtraction on this integer cost $O_H(p)$ bit
operations, not constant time on a machine with a fixed word width.
To read out, copy the index into scratch, recover its $m$ symbols,
and evaluate the $m$ affine maps. Per-step certified rounding error
$\epsilon_r$ contributes at most $\epsilon_r/(1-\lambda)$.
Choosing $O_H(p)$ fractional bits and a final output error at most
$2^{-p-2}$ leaves total numerical error at most $2^{-p-1}$.
The stated readout and storage costs follow for fixed-dimensional
affine maps and the stipulated output evaluator.

For completeness the separated homogeneous attractor is Ahlfors
regular of dimension $s=\log d/\log(1/\lambda)$: assign mass $d^{-m}$
to each depth-$m$ cylinder. A cylinder containing $x$ and of diameter
less than $r$ supplies the lower ball-mass bound. At a depth with
$\lambda^m$ comparable to $r$, distinct cylinders have separation
bounded below by a constant times $\lambda^m$; Euclidean packing
bounds the number intersecting a ball of radius $r$ by a model
constant, giving the upper bound. Reachable states are dense, since
the images of $x_0$ lie in every cylinder at every depth.
Theorem~\ref{III:M:thm:ahlfors-memory} supplies the matching lower bound.
\end{proof}

\begin{theorem}[An effective contractive coordinate chart]
\label{III:M:thm:chart-memory}
Suppose predictive states admit a supplied global chart
$\phi:[0,1]^k\to\mathcal X$ such that the conjugate updates
$F_a=\phi^{-1}\tau_a\phi$ map the cube to itself and are
$\lambda$-Lipschitz in maximum norm, $\lambda<1$.
Assume reachable states are dense and
\[
 c\|z-z'\|_\infty\le d_{\rm fut}(\phi z,\phi z')
                         \le C\|z-z'\|_\infty.
\]
Suppose the initial coordinate, the $F_a$ and $h\circ\phi$ have
certified $q$-bit evaluators using polynomial time, $O_H(q)$ scratch
bits and $O_H(q)$ immutable parameter bits, with fixed $k,d$.
Then an executable without a transition table attains error $2^{-p}$
with $kp+O_H(1)=B_f^*(2^{-p})+O_H(1)$ persistent bits and polynomial
compilation, update and readout bit cost.
\end{theorem}
\begin{proof}
Set $\epsilon=2^{-p}$ and choose a mesh $\delta=2^{-q}$ with
$\delta\le(1-\lambda)\epsilon/(4C)$, so $q=p+O_H(1)$.
Store the $k$ integer grid coordinates in $\{0,\ldots,2^q\}$.
Evaluate each update to maximum-norm error at most $\delta/2$,
clip to the cube and round to the nearest grid point. Clipping is
nonexpansive; the total one-step numerical error is at most $\delta$.
Initialize with error at most $\delta$. Thus
$e_{t+1}\le\lambda e_t+\delta$ and
$e_t\le\delta/(1-\lambda)$ for every history.
Output error from state tracking is at most $\epsilon/4$; evaluating
the readout to error at most $\epsilon/2$ meets the target.
Each coordinate needs $q+1$ bits, giving the asserted upper bound
and the stated computational costs by the supplied evaluators.
A grid packing of the cube at spacing greater than $2\epsilon/c$,
and reachable density, forces $\Omega_H(\epsilon^{-k})$ distinct
executable states by future observability. This gives the lower bound.
\end{proof}

The chart theorem assumes contraction in the executable coordinates;
bi-Lipschitz equivalence alone can enlarge a contraction constant
beyond one. Discovering such a chart is not part of the theorem.

Theorem~\ref{III:M:thm:packed-startup} below removes the fixed-point
initialization restriction for effective initial coordinates by encoding
startup lengths within the state index; it needs no address of $x_0$.

\paragraph{Why unequal ratios are a different target.}
For $\tau_0(x)=x/4$ and $\tau_1(x)=x/2+1/2$, the separated attractor
has dimension $s$ determined by $4^{-s}+2^{-s}=1$, hence
$s=\log_2((1+\sqrt5)/2)<1$.
A fixed-length register based on the worst contraction $1/2$ stores
$p+O(1)$ bits, rather than $sp+O(1)$. Thus the optimal
coefficient does not follow from that fixed-window construction for unequal ratios. Uniform bi-Lipschitz
bounds alone also do not prove Ahlfors regularity for arbitrary
nonconformal systems. Adaptive address coding or effective charts need separate theorems.
The rational-similarity case is supplied next in
Section~\ref{III:M:audit:sec:addresses}, with its own update and storage costs.

\subsection{Packed startup states remove the fixed-point initialization restriction}
\label{III:M:sec:packed-startup}
\begin{theorem}[Homogeneous addresses from an arbitrary effective initial state]
\label{III:M:thm:packed-startup}
Retain the separated homogeneous similarities and future-observability
assumptions of Theorem~\ref{III:M:thm:address-memory}, but replace its
fixed-point assumption by an arbitrary $x_0\in\mathcal X$ admitting
certified $q$-bit evaluation in polynomial time and $O_H(q)$ scratch
and immutable parameter bits. Assume also that $h$ has an effectively
evaluable $C_{\rm obs}$-Lipschitz extension to a fixed neighborhood
of $\mathcal X$, with the same evaluation resource bounds as before.
Then the same optimal persistent memory
\[
 B(p)=\frac{\log d}{\log(1/\lambda)}p+O_H(1)
       =B_f^*(2^{-p})+O_H(1)
\]
is achieved without a transition table or a supplied address of $x_0$.
Updates use $O_H(p)$ bit operations. Readout uses
$O_H(p\mathsf M(p))$ bit operations plus evaluation of $x_0$ and $h$,
and $O_H(p)$ scratch and immutable parameter bits.
\end{theorem}
\begin{proof}
Choose $m$ as in Theorem~\ref{III:M:thm:address-memory}, and retain the
entire input word until length $m$, and thereafter its last $m$
letters. The state set is $Q_m=\bigcup_{\ell=0}^m\Sigma^\ell$.
Its cardinality is $(d^{m+1}-1)/(d-1)$, so it uses
$m\log_2d+O_d(1)$ bits. Storing a separate length counter is
unnecessary: for a length-$\ell$ word of base-$d$ value $v$, put
\[
 S_\ell=\frac{d^\ell-1}{d-1},\qquad I=S_\ell+v.
\]
The length classes occupy the disjoint consecutive intervals
$[S_\ell,S_{\ell+1})$. Initialize $I=0$. On input
$a\in\{0,\ldots,d-1\}$, the exact update is
\[
 I\leftarrow
 \begin{cases}
 dI+a+1,&I<S_m,\\
 S_m+\big(d(I-S_m)+a\big)\bmod d^m,&I\ge S_m.
 \end{cases}
\]
This needs only the immutable $O_H(p)$-bit constants $S_m,d^m$
and fixed-$d$ integer arithmetic; hence the bit cost is $O_H(p)$.
No time-dependent external counter is used.

The representative of a stored word $u$ is $\tau_u(x_0)$, in the
sequential convention of Theorem~\ref{III:M:thm:address-memory}.
Before the window fills, this is the true state. Afterwards the
representative and the true state share their last $m$ contractions,
and so are within $\lambda^m\operatorname{diam}\mathcal X$.
The truncation output error is at most $2^{-p-1}$.

For readout, a scratch copy of $I$ decodes the letters backwards by
$a=(I-1)\bmod d$ and $I\leftarrow\lfloor(I-1)/d\rfloor$
until zero; a scratch word stores the reversed letters for replay.
There are at most $m=O_H(p)$ such operations, each on $O_H(p)$ bits.
Evaluate $x_0$ and replay the affine maps with initial error
$e_0\le 2^{-p}/(8C_{\rm obs})$ and per-step error
$e_r\le(1-\lambda)2^{-p}/(8C_{\rm obs})$.
The total state error is at most
$e_0+e_r/(1-\lambda)\le2^{-p}/(4C_{\rm obs})$.
A final readout error at most $2^{-p-2}$ completes the budget.
Fixed model constants ensure the computation stays in the stipulated
neighborhood; finitely many small precisions can use finer arithmetic.
The stated storage and readout bounds follow.
For every $x_0\in\mathcal X$, its finite images meet every cylinder,
so reachability is dense. The same packing lower bound as in
Theorem~\ref{III:M:thm:address-memory} applies.
\end{proof}

This removes an initialization restriction within the homogeneous
class. Section~\ref{III:M:audit:sec:addresses} now treats unequal rational
ratios; discovering effective coordinates for a general regular
attractor remains open.

\section{Unequal rational similarity addresses}\label{III:M:audit:sec:addresses}
This section gives a constructive answer to the
unequal-rational-ratio part of the effective-memory target. Its ingredients
include classical stopping trees, enumerative coding, and self-similar packing.
In particular, Reznik--Anisimov~\cite{reznik} already use composition groups and
combinatorial indexing to avoid complete variable-to-fixed coding trees.
Those coding ideas are not claimed as new here. The proposed addition is their
explicit use in an autonomous predictive-state update, with initialization,
finite precision, and bit costs stated together.

\begin{theorem}[Streaming rational-similarity addresses]
\label{III:M:audit:thm:addresses}
Fix a finite alphabet of size $d\ge2$, a dimension $k$, and rational affine
similarities
\[
 \tau_a(x)=r_aO_ax+t_a,\qquad O_a^\top O_a=I,\qquad
 r_a\in\mathbb Q\cap(0,1).
\]
Assume their compact attractor $X$ has positive diameter and strongly separated
first-level images. The affine coefficients and all model constants are fixed.
Let $x_0\in X$ be effectively evaluable. Suppose the scalar output $h$ has a
certified $C$-Lipschitz extension to a fixed neighborhood of $X$, and
\[
 c\Mnorm{x-y}\le d_{\mathrm{fut}}(x,y)
 :=\sup_{v\in\Sigma^*}|h(\tau_vx)-h(\tau_vy)|
 \le C\Mnorm{x-y},\qquad c,C>0.
\]
Suppose $x_0$ and $h$ have certified $q$-bit evaluators with polynomial time
and $O_H(q)$ scratch and immutable parameter bits. Here $H$ denotes the fixed
model, not an input whose size is being suppressed in a uniform claim.
Let $s>0$ be the unique solution of $\sum_a r_a^s=1$.
Then, for error $2^{-p}$, an autonomous finite-state executable attains
\[
 B(p)=sp+O_H(1)=B_f^*(2^{-p})+O_H(1)
\]
persistent bits, polynomial compilation/update/readout bit cost, and $O_H(p)$
scratch and immutable parameter bits. No exponential word or transition table
is stored. A conservative schoolbook implementation of the indexing operations
uses $O_H(p^{d+4})$ bit operations per update or readout, in addition to the
stipulated evaluator cost. This is not a uniform low-degree bound as $d$ varies.
\end{theorem}

\begin{proof}
\emph{Stopping address and update.}
Take an integer constant $c_0\ge1$ with $C\operatorname{diam}(X)2^{-c_0}
\le1/2$, and set $\theta=2^{-(p+c_0)}$. Write a history backwards, newest
letter first. Keep its shortest prefix $v$ with contraction product
$r_v=\prod_i r_{v_i}\le\theta$, or the entire reversed history if the
threshold has not been reached. The possible addresses are all internal nodes
and leaves of the finite ordered tree which stops a branch as soon as its
product is at most $\theta$.

On a new input $a$, prepend $a$ and apply this stopping rule again.
This uses only the old address and $a$. If the old address was already a leaf,
the new stopping prefix ends within that address; discarded older letters
are never needed. If it was internal, the full old history was retained.
This proves the update invariant and includes startup without a separate clock.
The representative replays the retained letters chronologically from $x_0$.
Before truncation it is exact. Thereafter the true state and representative
share all retained contractions, so
\[
 \Mnorm{x_t-\widehat x_t}\le r_v\operatorname{diam}(X)
 \le\theta\operatorname{diam}(X).
\]
The exact readout error is at most $2^{-p-1}$.

\emph{Number of addresses.}
Put $q_a=r_a^s$, so $\sum_a q_a=1$, and let $L_\theta$ be the number of
leaves. The stopped tree is finite and complete. Its leaf masses sum to one,
and, with $r_{\min}=\min_a r_a$,
\[
 r_{\min}\theta<r_v\le\theta
 \quad\Longrightarrow\quad
 (r_{\min}\theta)^s<q_v\le\theta^s.
\]
Hence $\theta^{-s}\le L_\theta<(r_{\min}\theta)^{-s}$. A full $d$-ary
tree with $L_\theta$ leaves has
\[
 N_\theta=\frac{dL_\theta-1}{d-1}
\]
total nodes, including all startup states. An ordinal in
$\{0,\ldots,N_\theta-1\}$ thus needs $sp+O_H(1)$ bits.

\emph{Counting without listing words.}
For a node of contraction product $R>\theta$, define
\[
 I(R)=\sum_{\substack{n\in\MN^d\\MR\prod_a r_a^{n_a}>\theta}}
       \binom{|n|}{n_1,\ldots,n_d},
 \qquad N(R)=1+dI(R).
 \tag{M.1}
\]
For $R\le\theta$, set $N(R)=1$. Every word whose product is above the
threshold is internal, since all shorter prefixes have a larger product.
The multinomial counts all such ordered words of a composition $n$ exactly.
Thus $I(R)$ is the number of internal nodes below this node; the full-tree
identity gives $N(R)$. Each coordinate of $n$ is $O_H(p)$ and there are
$O_H(p^d)$ compositions. This groups words for \emph{counting only}; distinct
letter orders retain distinct ordinals and distinct map compositions.

Use preorder indexing. At each internal node its empty continuation has
index zero, and its ordered child subtrees have sizes $N(Rr_a)$.
Subtracting or adding these sizes performs rank/unrank. There are $O_H(p)$
letters in every address. All threshold comparisons are exact rational
comparisons; neither $s$ nor a floating-point logarithm is needed by the codec.

For completeness, enumerate composition counters one at a time, compute
products and multinomials afresh, and retain only the sum. A multinomial can
be computed as a product of binomial coefficients using multiplicative exact
recurrences. All retained arithmetic values have $O_H(p)$ bits; full factorials
are unnecessary. Each composition costs at most $O_H(p^3)$ schoolbook bit
operations. Rank/unrank makes $O_H(p)$ counting calls, giving the stated
conservative time and linear scratch bounds. This deliberately trades time
for space; a faster implementation may use a separately charged cache.

\emph{Finite-precision readout.}
Replay $O_H(p)$ rational affine maps. With $r_{\max}<1$, initial and per-step
errors satisfying
\[
 e_0\le\frac{2^{-p}}{8C},\qquad
 e_r\le\frac{(1-r_{\max})2^{-p}}{8C}
\]
give state error at most $2^{-p}/(4C)$. A final readout error at most
$2^{-p-2}$ meets the total budget. The fixed neighborhood assumption makes
the rounded readout well defined. Finitely many small precisions may use
finer arithmetic with a fixed model-dependent overhead.

\emph{Lower bound.}
Strongly separated self-similar attractors are $s$-Ahlfors regular, also when
the ratios differ. One elementary proof assigns cylinder mass $r_v^s$ and
uses stopping cylinders with diameters comparable to a ball radius: separation
bounds how many meet the ball, and a cylinder through its center gives the
lower mass bound. The orbit of any $x_0\in X$ is dense, because it meets every
cylinder. Future observability and the packing argument of Section~\ref{III:M:sec:memory},
therefore force $\Omega_H(2^{sp})$ distinguishable executable states.
This gives $B_f^*(2^{-p})\ge sp-O_H(1)$ and completes the proof.
\end{proof}


\textbf{Scope.} Rational ratios and fixed alphabet size make the exact
polynomial counting argument effective. This proof supplies neither efficient
discovery of the model nor the general nonconformal or homographic state encoder.
Persistent bits are optimal to additive constant; the leading coefficient of
\emph{total} peak storage has not been optimized. The generic implementation
is a proof of polynomial cost, not a latency improvement.

\subsection{The unequal-ratio example already singled out in the preceding development}
For
\[
 \tau_0(x)=x/4,\qquad\tau_1(x)=x/2+1/2,\qquad h(x)=x,
\]
the attractor has diameter one, the first-level images are separated, and
$d_{\rm fut}(x,y)=|x-y|$. Its dimension solves
\[
 4^{-s}+2^{-s}=1,\qquad s=\log_2\frac{1+\sqrt5}{2}
 \simeq0.694241914.
\]
At threshold $\theta=2^{-\ell}$, the stopped tree has $F_{\ell+2}$ leaves
and $2F_{\ell+2}-1$ nodes, where $F_0=0,F_1=1$. The leading persistent-bit
coefficient is about 30.6\% below the worst-contraction fixed binary window.
This percentage concerns a mathematical state representation; no measured
memory or runtime improvement for a deployed system is asserted.

\begin{center}
\begin{tabular}{@{}rrr@{}}\toprule
Threshold exponent $\ell$ & Stopping-address bits & Fixed-window bits, with startup\\\midrule
8 & 7 & 9\\16 & 13 & 17\\32 & 24 & 33\\64 & 46 & 65\\128 & 91 & 129\\\bottomrule
\end{tabular}
\end{center}
These are exact finite address counts. A finite-precision output budget needs
the constant tolerance allocation used in the theorem.

\section{An exact finite optimum for the Fibonacci memory}
\label{III:M:sec:leaf-optimum}
The 91-bit entry in the preceding table counts every node of the stopped
tree, including startup nodes. It is a correct cost for that construction.
This section changes the startup contract: the prescribed approximation
error is allowed from time zero. For the explicit initial state $x_0=0$,
leaf addresses alone then give an exact finite minimax result. No claim
about an arbitrary initial point, exact startup, or total peak bytes is
implicit in this improvement.

Fix an integer $\ell\ge1$, put $\theta=2^{-\ell}$, and consider
\[
 x_{t+1}=\tau_{a_{t+1}}(x_t),\qquad
 \tau_0(x)=x/4,\quad \tau_1(x)=x/2+1/2,\quad x_0=0.
\]
The readout is $h(x)=x$. A deterministic memory consists of a finite set
$Q$, a fixed initial state, updates $U_a:Q\to Q$ and a readout
$D:Q\to\MR$, all independent of an uncharged external clock or transcript.
Its required error is $|D(q_t)-x_t|\le\theta/2$ for every finite input
history, including the empty history. Persistent model constants and the
fixed precision parameter are treated as in the earlier fixed-model
contract; temporary arithmetic is accounted for separately below.

Write addresses from newest to oldest. For a word $v=v_1\cdots v_n$, let
\[
 \tau_v=\tau_{v_1}\circ\cdots\circ\tau_{v_n},\qquad
 c(v)=2\#_0(v)+\#_1(v),\qquad I_v=\tau_v([0,1]).
\]
Let $\mathcal L_\ell$ be the words whose cost first reaches or exceeds
$\ell$: $c(v)\ge\ell$ and $c(v^-)<\ell$, where $v^-$ deletes the last
letter. These are leaves, not all prefixes of leaves.

\begin{lemma}[Exact leaf count and separated endpoints]
\label{III:M:lem:fibonacci-leaf-packing}
With $F_0=0,F_1=1$, one has $|\mathcal L_\ell|=F_{\ell+2}$.
Every leaf interval has length $\theta$ or $\theta/2$.
The right endpoints $b_v=\tau_v(1)$, ordered from left to right, have
successive separation at least $3\theta/2$.
Consequently the minimum number of closed intervals of length $\theta$
covering the attractor $K$ is exactly $F_{\ell+2}$.
\end{lemma}
\begin{proof}
The leaf count $L(n)$ satisfies $L(n)=L(n-2)+L(n-1)$ for $n>0$ and
$L(n)=1$ for $n\le0$. Thus $L(\ell)=F_{\ell+2}$. The last increment
is either one or two, so a leaf has cost $\ell$ or $\ell+1$.

Consider adjacent leaves and their longest common prefix $u$, of
contraction ratio $r_u$. The left leaf is the rightmost descendant of
$u0$, so its right endpoint is $\tau_u(1/4)$, because $\tau_1(1)=1$.
The right leaf is the leftmost descendant of $u1$, with some number of
zeros afterwards. Its left endpoint is $\tau_u(1/2)$ and its length is
some $r_v\in\{\theta,\theta/2\}$. The right-endpoint separation is
$r_u/4+r_v$. Since $u$ is internal, its dyadic ratio satisfies
$r_u\ge2\theta$. If $r_u=2\theta$, the right child is already a leaf
and $r_v=\theta$. Otherwise $r_u\ge4\theta$ and $r_v\ge\theta/2$.
Both cases give the asserted $3\theta/2$ bound.
Each endpoint belongs to $K$. A length-$\theta$ interval contains at
most one of these endpoints, while the $F_{\ell+2}$ leaf intervals
themselves cover $K$ and have length at most $\theta$.
\end{proof}


\begin{theorem}[Exact autonomous optimum with approximate startup]
\label{III:M:thm:fibonacci-exact-optimum}
Under the stated contract and $x_0=0$, the minimum number of executable
memory states and the minimum fixed-length persistent bit budget are
\[
 |Q|_{\min}=F_{\ell+2},\qquad
 B_{\min}(2^{-\ell-1})=\left\lceil\log_2 F_{\ell+2}\right\rceil.
\]
In particular, $B_{\min}(2^{-129})=90$ bits at $\ell=128$.
\end{theorem}
\begin{proof}
For the upper bound take $Q=\mathcal L_\ell$, initialize at
$v_0=0^{\lceil\ell/2\rceil}$, update by
\[
 U_a(v)=\operatorname{trim}_\ell(av),\qquad
 D(v)=\tau_v(1/2),
\]
where trimming retains the shortest prefix of cost at least $\ell$.
Such a prefix always exists in $av$. Initially $x_0\in I_{v_0}$.
If $x_t\in I_v$ and $av=v'u$, invariance of $[0,1]$ implies
$x_{t+1}\in\tau_{v'}(\tau_u([0,1]))\subseteq I_{v'}$.
Induction proves this inclusion for all histories. The midpoint readout
therefore has error at most $|I_{v'}|/2\le\theta/2$, including startup.
Every leaf is reachable by supplying its letters in chronological order,
since $\operatorname{trim}_\ell(vv_0)=v$.

For the lower bound one must use reachable histories, rather than assume
that every point of $K$ occurs at finite time. Set $T=\ell+3$.
For each leaf $v$ first supply $T-|v|$ copies of letter 1 and then the
letters of $v$ in reverse order. The final true state is
\[
 y_v=\tau_v\bigl(\tau_1^{T-|v|}(0)\bigr),\qquad
 |b_v-y_v|=r_v2^{-(T-|v|)}\le\theta/8,
\]
because $|v|\le\ell$. Hence two different $y_v$ are separated by at
least $3\theta/2-2\theta/8=5\theta/4>\theta$.
All these histories have the same length. If two ended in the same
stored state, their common readout would be within $\theta/2$ of both,
contradicting this separation. Thus at least $F_{\ell+2}$ states are
necessary. Finally
$F_{130}=659034621587630041982498215$ lies strictly between $2^{89}$
and $2^{90}$, proving the numerical bit optimum.
\end{proof}


\begin{proposition}[Leaf identifiers without an exponential table]
\label{III:M:prop:fibonacci-leaf-cost}
The preceding memory can store only a leaf ordinal. With elementary
integer arithmetic its update and exact midpoint readout use $O(\ell)$
temporary bits and at most $O(\ell^3)$ bit operations per update.
No list of the $F_{\ell+2}$ leaves or transition table is required.
\end{proposition}
\begin{proof}
At remaining cost $n>0$, the left subtree contains $L(n-2)$ leaves.
For ranking, add this number when the next letter is 1, then decrease
$n$ by the letter cost. For unranking, compare the current ordinal with
$L(n-2)$ and choose the corresponding child, subtracting that count on
the right. There are at most $\ell$ letters. Each needed Fibonacci
number can be recomputed by a rolling two-integer recurrence, with
$O(\ell)$ additions of $O(\ell)$-bit integers. Recomputing rather than
caching all counts gives $O(\ell^3)$ bit operations and $O(\ell)$ live
bits for the ordinal, counters and retained word. Prepend, trim and rank
give the autonomous update. The midpoint is a dyadic rational with at
most $\ell+2$ fractional bits, obtainable by exact affine replay within
the same space and a smaller time bound.
\end{proof}


\paragraph{What changes, and what does not.}
The earlier $2F_{130}-1$-node representation still uses 91 bits and
retains exact startup histories. The new 90-bit representation uses a
midpoint readout and permits the declared error immediately. The
91-bit count was never a lattice quantization of occupation parameters.
The asymptotic coefficient remains
$\log_2\varphi\simeq0.694241914$, with
$\varphi=(1+\sqrt5)/2$; improving startup changes an additive constant.
Hausdorff dimension alone would only give a rate with unspecified
finite constants. The exact packing above supplies the missing finite
argument. Independently byte-aligned 90-bit and 91-bit payloads both
occupy 12 bytes; saving one bit is not automatically a resident-byte gain.
Python object sizes, working storage, output buffers and time are not
included in the persistent-bit number.


\subsection{Extension to the leaf union and non-dyadic tolerances}\label{III:M:rev13:sec:fib}

Notation as in Section~\ref{III:M:sec:leaf-optimum}: $\tau_0(x)=x/4$, $\tau_1(x)=x/2+1/2$, $h(x)=x$, $\theta=2^{-\ell}$, leaves
$\mathcal L_\ell$ (newest-first words whose cost first reaches $\ell$),
$I_v=\tau_v([0,1])$, $b_v=\tau_v(1)$, and $K$ the attractor. Put
\[
U_\ell=\bigcup_{v\in\mathcal L_\ell}I_v\supseteq K .
\]
The preceding theorem proves $|Q|_{\min}=F_{\ell+2}$ for $x_0=0$ and states that no claim for an
arbitrary initial point is implicit. The proof in fact needs only $x_0\in U_\ell$.

\begin{theorem}[Exact optimum for every initial point of the leaf union]
\label{III:M:rev13:thm:fib-general}
Fix $\ell\ge1$ and the contract $|D(q_t)-x_t|\le\theta/2$ for every finite history
including the empty one, for an autonomous deterministic memory $(Q,q_0,(U_a),D)$ without
an uncharged clock. For every $x_0\in U_\ell$ (in particular every $x_0\in K$),
\[
|Q|_{\min}=F_{\ell+2},\qquad B_{\min}(2^{-\ell-1})=\lceil\log_2F_{\ell+2}\rceil .
\]
For every $x_0\in[0,1]$ the lower bound $|Q|_{\min}\ge F_{\ell+2}$ holds. An exact state count for $x_0\notin U_\ell$ is not asserted here.
\end{theorem}
\begin{proof}
\emph{Upper bound.} Choose (or supply to the compiler) $v_0\in\mathcal L_\ell$ with $x_0\in I_{v_0}$; this is the only
place where $x_0$ enters. Take $Q=\mathcal L_\ell$, $U_a(v)=\Mtrim_\ell(av)$,
$D(v)=\tau_v(1/2)$. If $x_t\in I_v$ and $av=v'u$ with $v'=\Mtrim_\ell(av)$, then
$x_{t+1}=\tau_a(x_t)\in\tau_{av}([0,1])=\tau_{v'}(\tau_u([0,1]))\subseteq I_{v'}$, so by
induction $x_t\in I_{v_t}$ for all histories and the midpoint readout has error at most
$|I_{v_t}|/2\le\theta/2$.

\emph{Lower bound.} Let $T=\ell+3$. For each leaf $v$ feed $T-|v|$ letters $1$ followed by
the letters of $v$ in chronological order (possible because $|v|\le\ell$). The final state
is $y_v=\tau_v\bigl(\tau_1^{T-|v|}(x_0)\bigr)$ and, since $\tau_1^j(x_0)=1-2^{-j}(1-x_0)$,
\[
|b_v-y_v|=r_v2^{-(T-|v|)}|1-x_0|\le\theta\,2^{-3}=\theta/8
\]
for every $x_0\in[0,1]$. The right endpoints $b_v$ are pairwise at least $3\theta/2$ apart
(Lemma~\ref{III:M:lem:fibonacci-leaf-packing}), hence distinct $y_v$ are
at least $3\theta/2-\theta/4=5\theta/4>\theta$ apart. All $F_{\ell+2}$ histories have the
same length $T$; two of them ending in the same memory state would force a common
readout within $\theta/2$ of both final states, contradicting the separation.
\end{proof}


\begin{remark}
The leaf union $U_\ell$ is a strict superset of $K$ (it contains the gaps of $K$ at
depth $\ge\ell$), so the theorem covers more initial points than the attractor. For
$x_0\in[0,1]\setminus U_\ell$ (for example $x_0=3/10$ at $\ell\ge2$) the leaf machine
cannot be initialised, the exact value lies in $[F_{\ell+2},\,2F_{\ell+2}-1]$ (the
all-node construction gives the upper end). The dependence of the precise
count on $x_0$ in this residual case remains open.
\end{remark}

\begin{proposition}[Non-dyadic tolerances]\label{III:M:rev13:prop:fib-eps}
Let $x_0\in K$, let $0<\epsilon\le1/4$ and let $\ell\ge1$ be the integer with
$2^{-\ell-1}\le\epsilon<2^{-\ell}$, $\theta=2^{-\ell}$. For the contract
$|D(q_t)-x_t|\le\epsilon$:
\[
\text{if }\ \tfrac\theta2\le\epsilon<\tfrac{3\theta}4:\quad |Q|_{\min}(\epsilon)=F_{\ell+2};
\qquad
\text{if }\ \tfrac{3\theta}4\le\epsilon<\theta:\quad F_{\ell+1}\le|Q|_{\min}(\epsilon)\le F_{\ell+2}.
\]
\end{proposition}
\begin{proof}
The leaf machine at threshold $\theta\le2\epsilon$ has error $\le\theta/2\le\epsilon$, giving
the upper bound. For the lower bound repeat the construction of
Theorem~\ref{III:M:rev13:thm:fib-general} with $T=\ell'+j$ at threshold $\theta'=2^{-\ell'}$: the
final states are within $\theta'2^{-j}$ of the endpoints, hence pairwise more than
$3\theta'/2-\theta'2^{1-j}$ apart, and this exceeds $2\epsilon$ for large $j$ whenever
$\epsilon<3\theta'/4$. Take $\ell'=\ell$ in the first case and $\ell'=\ell-1$ (so
$3\theta'/4=3\theta/2>\epsilon$) in the second.
\end{proof}


\section{Persistent facts, checkpoint replay and certified guards}\label{III:M:audit:sec:persistent}
The new address construction still forgets old distinctions under a sufficiently
long common continuation. It does not solve permanent factual recall.

\begin{proposition}[Uniform contraction cannot preserve a non-fading fact]
Suppose two histories encode a binary fact $b\in\{0,1\}$ that remains unchanged
after $t$ repetitions of an irrelevant event. The two stored states initially
have distance at most $D$. Suppose that event is $\lambda$-contractive,
$0<\lambda<1$, and the answer distribution is $L$-Lipschitz in state distance
under total variation. If both fact answers have TV error at most $\epsilon$
against the correct point-mass answers, then
\[
 1\le 2\epsilon+LD\lambda^t.
\]
Consequently no fixed $\epsilon<1/2$ can hold for all $t$ under these assumptions.
\end{proposition}
\begin{proof}
The correct answers have TV distance one. The triangle inequality bounds this
by the two answer errors plus the distance between the model's two answer
distributions. Lipschitz readout and contraction bound the latter by
$LD\lambda^t$.
\end{proof}

This is an elementary obstruction, not a new novelty claim. Discontinuous
exact-real readout is excluded by the Lipschitz hypothesis; finite robust
storage cannot silently use infinitely precise coordinates.

If all $N$ independently supplied bits must be individually recalled with
worst-case error below $1/2$ by a deterministic-state encoder, at least $N$
persistent bits are necessary. Otherwise two of the $2^N$ fact assignments
share a state and disagree on a query. Thus growing factual content cannot be
hidden inside a fixed total byte budget with unrestricted recall guarantees.

\subsection{The persistent-plus-predictive construction target}
\textbf{Target: a certified streaming memory with exact persistent distinctions
and a compressed predictive component.} The deployed state has the form
\[
 m_t=(r_t,z_t),\qquad
 r_{t+1}=G(r_t,a_{t+1}),\qquad z_{t+1}=F_{a_{t+1},r_t}(z_t).
\]
Here $r_t$ retains designated non-fading facts or exceptions, while $z_t$
captures compressible dynamics. Exact storage refers to the chosen typed
facts, not to infallible extraction from natural language. Every byte in
$r_t$, provenance and indexes is counted. A guarantee must also account for
coupling from $z_t$ into any update of $r_t$.

\subsection{Constructive extension: replay from the mode at the window base}
\begin{proposition}[Rational addresses with a mode checkpoint]
\label{III:M:audit:prop:checkpoint}
Let $\mathcal R$ be a nonempty finite set, $G:\mathcal R\times\Sigma\to
\mathcal R$ a deterministic computable transition, and
\[
 r_{t+1}=G(r_t,a_{t+1}),\qquad
 x_{t+1}=\tau_{a_{t+1},r_t}(x_t),\qquad
 \tau_{a,r}(x)=r_aO_{a,r}x+t_{a,r}.
\]
The ratios $r_a\in\mathbb Q\cap(0,1)$ depend only on the letter; the
orthogonal maps and translations may depend on the mode. Assume a common
nonempty compact invariant set $X$ contains $x_0$, and let $D$ be a known
upper bound on its diameter. Fix $0<\theta<1$ and use the stopped address
tree of Theorem~\ref{III:M:audit:thm:addresses}, including its internal startup nodes.
An autonomous representation is
\[
 m_t=(b_t,I_t),\qquad b_t=r_{\mathrm{base},t}
       =r_{t-|v_t|},\qquad I_t=\operatorname{rank}(v_t),
\]
where $v_t$ is the retained word, newest letter first. It uses at most
\[
 \lceil\log_2|\mathcal R|\rceil+\lceil\log_2N_\theta\rceil
\]
persistent bits. Start with $(b_0,I_0)=(r_0,0)$. On receiving $a$, form
\[
 w=(a)\operatorname{unrank}(I_t)=v_{t+1}u_{\rm drop},\qquad
 v_{t+1}=\operatorname{trim}(w),
\]
and update, with $G^*$ iterating $G$ in the displayed word order,
\[
 b_{t+1}=G^*(b_t,\operatorname{rev}(u_{\rm drop})),\qquad
 I_{t+1}=\operatorname{rank}(v_{t+1}).                 \tag{M.2}
\]
To read, initialize $(\rho,y)=(b_t,x_0)$ and for each letter $a$ of
$\operatorname{rev}(v_t)$ apply the simultaneous assignment
\[
 (\rho,y)\longleftarrow\bigl(G(\rho,a),\tau_{a,\rho}(y)\bigr),
                                                        \tag{M.3}
\]
whose right-hand side uses the old mode. The final mode is exactly $r_t$.
If all history is retained, $y=x_t$; otherwise
\[
 \Mnorm{x_t-y}\le D r_{v_t}\le D\theta.               \tag{M.4}
\]
Neither invertibility nor strong connectivity of $G$ is required.

For executable numerical bounds, assume additionally rational affine
coefficients and the certified evaluation hypotheses of
Theorem~\ref{III:M:audit:thm:addresses}, uniformly over the fixed finite mode catalogue.
Then the same polynomial indexing cost and $O_H(p)$ scratch bounds apply
at $\theta=2^{-(p+c_0)}$. For an output uniformly $C$-Lipschitz in the
continuous coordinate within each exact mode, the same error allocation
gives output error $2^{-p}$ and
\[
 B(p)\le sp+\lceil\log_2|\mathcal R|\rceil+O_H(1),
 \qquad \sum_a r_a^s=1.                              \tag{M.5}
\]
All mode/transition descriptions and numerical evaluators are charged.
This is a constructive upper bound, not a joint optimality theorem.
\end{proposition}
\begin{proof}
Let $\ell=|v_t|$. Its oldest endpoint is time $t-\ell$, whose mode is $b_t$.
After appending one letter, discarding $j=|u_{\rm drop}|$ oldest letters
advances that endpoint to $t-\ell+j=t+1-|v_{t+1}|$. Those discarded letters
are precisely $\operatorname{rev}(u_{\rm drop})$ in chronological order.
Equation (M.2) therefore advances the checkpoint to the correct new mode,
proving its invariant by induction. This only evaluates $G$ forwards.

Starting at this checkpoint, (M.3) reconstructs every subsequent mode exactly,
because $G$ depends on that mode and the input letter alone. The true state
at the window base and $x_0$ both lie in $X$. The two replays apply identical
mode-dependent similarities, whose composition has contraction ratio
$r_{v_t}$. This proves (M.4). If the whole history is retained, its base state
is $x_0$ and the error is zero. An internal tree node necessarily retains the
whole history. A leaf can also occur at the first threshold crossing before
any deletion; (M.4) is still a safe upper bound there.

The ratio depends only on the letter, so the address tree and composition
counts are exactly those of Theorem~\ref{III:M:audit:thm:addresses}. The one extra mode
uses $\lceil\log_2|\mathcal R|\rceil$ bits. Replaying or dropping letters
adds only $O_H(p)$ mode transitions. With certified rounded affine replay,
the numerical error satisfies the same geometric recurrence as before;
the discrete mode remains exact. This proves the computational claims.
\end{proof}


Strong separation and future observability are not needed for this upper
bound. They cannot be silently transferred from the unmixed family to the
mode-dependent one to claim a matching lower bound. Some pairs $(b,v)$ can
be unreachable, and different modes can have indistinguishable futures.
The model is fixed in (M.5); a growing mode catalogue must be charged explicitly.

\textbf{Feedback boundary.} If instead
$r_{t+1}=G(r_t,a_{t+1},\sigma_{t+1})$, where $\sigma_{t+1}$ is determined
by the continuous trajectory, the letters and checkpoint need not recover
the branch choices. One must prove that each choice is reconstructible,
or retain the branch information and charge its storage and update costs.
A log of one independent binary choice per retained step costs
$O_H(p)$ additional bits; no additive-constant overhead follows in general.
Proposition~\ref{III:M:audit:prop:checkpoint} does not resolve this feedback case.

\subsection{Numerical refinement, geometric envelopes and actual ambiguity}
Fix a stored state $m=(b,v)$, a query/event, and its scalar guard $g$.
Let $\Phi_{b,\operatorname{rev}(v)}$ denote chronological replay with modes
initialized at $b$. For a leaf address put
\[
 c_m=\Phi_{b,\operatorname{rev}(v)}(x_0),\qquad
 S_m=\Phi_{b,\operatorname{rev}(v)}(X),\qquad R_m=Dr_v.
\]
The set of states attainable by admissible histories with this \emph{full}
stored state is
\[
 \mathcal C_m=\{x(h):\operatorname{Enc}(h)=m\}
       \ \subseteq\ S_m\ \subseteq\ B(c_m,R_m).       \tag{M.6}
\]
In general the first inclusion is strict: arbitrary points of $X$ need not
be reachable with the required base mode and encoding. For an internal
startup node use $S_m=\{c_m\}$ and $R_m=0$ instead. All these histories
have the same current mode, which is recovered by (M.3).

\begin{lemma}[Termination of a centered Lipschitz sign certificate]
\label{III:M:audit:lem:guard}
Suppose $g$ is $L_g$-Lipschitz on a region containing the relevant states and
evaluation points. At precision $q$ let $\widetilde g_q$ be a certified
evaluation of the exact replay center with
\[
 |\widetilde g_q-g(c_m)|\le\delta_q,
 \qquad \delta_q=L_gE_{\rm num}(q)+\eta_g(q)\longrightarrow0.
\]
Put $A=L_gR_m$ and $\Delta_m=|g(c_m)|-A$. The test
\[
 |\widetilde g_q|>A+\delta_q                          \tag{M.7}
\]
certifies a constant strict sign on $\mathcal C_m$. Eventual acceptance
by this particular test is equivalent to $\Delta_m>0$; a sufficient
precision is $2\delta_q<\Delta_m$.
The alternative test
\[
 |\widetilde g_q|+\delta_q\le A                      \tag{M.8}
\]
certifies $\Delta_m\le0$, so further numerical refinement cannot make
(M.7) succeed. If $\Delta_m<0$, (M.8) eventually succeeds. If $\Delta_m=0$,
neither finite-time exit is guaranteed.
\end{lemma}
\begin{proof}
For any $x\in\mathcal C_m$, $|g(x)-\widetilde g_q|\le A+\delta_q$,
which proves the sign claim. Acceptance implies
$|g(c_m)|\ge|\widetilde g_q|-\delta_q>A$. Conversely,
$|\widetilde g_q|\ge|g(c_m)|-\delta_q>A+\delta_q$ whenever
$2\delta_q<\Delta_m$. The upper triangle inequality proves (M.8);
if $\Delta_m<0$, its left side is at most $|g(c_m)|+2\delta_q<A$
eventually. For $g(c_m)=A$ and
$\widetilde g_q=A+\delta_q/2$ with $\delta_q>0$, both (M.7) and (M.8)
fail for every $q$, even though the errors tend to zero.
\end{proof}


\textbf{Counterexample: numerical refinement is not a truncation certificate.}
Take $L_g=1$, $R_m=4/256$, $\delta_q=2/256$,
$\widetilde g_q=5/256$, and $g(c_m)=3/256$.
Then $|\widetilde g_q|>L_gR_m$, but $\Delta_m=-1/256$.
The weaker inequality omitting numerical uncertainty is invalid.
For instance, $c_m=0$, $g(x)=3/256+x$ on
$[-4/256,4/256]$ realizes a sign change despite that weaker test.

\textbf{Interval certificate.} Any certified enclosure
$[\underline g_m,\overline g_m]\supseteq g(\mathcal C_m)$ yields a
positive decision when $\underline g_m>0$, and a negative decision when
$\overline g_m<0$. An enclosure of $g(S_m)$ is sufficient.
Zero in the enclosure is inconclusive. Likewise $\Delta_m\le0$ only
excludes the centered test (M.7); it does not establish a sign change.
For $g(x)=R^2+x^2$ on $[-R,R]$, $R>0$, one has $L_g=2R$ and
$g(0)=R^2<2R^2=L_gR$, whereas the exact range
$[R^2,2R^2]$ is strictly positive.

\textbf{An actual information-loss certificate.} Exhibit two admissible
histories $h^+,h^-$ with identical \emph{full} stored state, the same
query/event, and $g(x(h^+))>0>g(x(h^-))$. No deterministic decision from
that state alone can be correct for both. This is stronger than finding
zero in an outer interval. A concrete example uses a single mode,
$\tau_0(x)=x/2$, $\tau_1(x)=x/2+1/2$, $x_0=0$, and threshold $1/2$.
Chronological histories $(0,0)$ and $(1,0)$ retain the same last-letter
address $(0)$, but their states are $0$ and $1/4$; the guard $x-1/8$
has opposite signs. Equality with zero requires a separately declared
tie convention and is not used in this example.

\textbf{Branching is a set-valued guarantee.} With feedback, unresolved
guards can be handled by keeping joint hypotheses about the discrete mode,
continuous compatible set and required replay/branch payload. Intersect each
compatible set with its guard region, then propagate it by that branch's
transition. This preserves containment of the reference trajectory when the
enclosures are sound. It does not provide a unique exact mode. A subset of
the original mode--address product alone need not encode this information,
because different branch histories can share its elements. Explicit binary
branching can double the number of hypotheses per unresolved step; merging
or symbolic compression requires a separate correctness and cost argument.

\subsection{Persistence is an observability property, not a spectral shortcut}
For $F(x)=x/(1+x)$ on $[0,1]$, direct induction gives
\[
 F^t(1)=\frac1{t+1}\longrightarrow0,\qquad
 \lim_{t\to\infty}(F^t(1))^{1/t}=1,\qquad F'(0)=1.
\]
With identity readout and this single neutral transition,
$d_{\rm fut}(F^t(1),0)=1/(t+1)$. Thus unit asymptotic exponential rate
does not imply a positive surviving separation.

For $P_0=\operatorname{diag}(1,0)$ and
$P_1=\operatorname{diag}(0,1)$, the joint spectral radius is one:
all products have norm at most one and repeated $P_0$ has norm one.
Nevertheless $P_1P_0=P_0P_1=0$, so a mixed continuation can erase every
distinction. These projectors \emph{commute}; noncommutativity is not needed
for the counterexample. Joint spectral radius controls the largest product
growth, not a lower bound for every continuation.

A word-Hankel operator connects pasts to future observables. Its nonzero
singular values count linear predictive rank, and their magnitudes measure
spectral approximation strength. They do not by themselves provide a
common center--stable decomposition for all letter transitions.

One \emph{sufficient specification} for finitely many protected facts is a
finite history congruence $\pi:\Sigma^*\to\mathcal R$ satisfying
$\pi(ha)=G(\pi(h),a)$, together with a designated fact label $f(r)$ and
a declared neutral alphabet. Require fact invariance under neutral events
and the uniform observational margin
\[
 \inf_{w\in\Sigma_{\rm neutral}^*}
 d_{\rm fut}(hw,h'w)\ge2\delta>0
 \quad\text{whenever } f(\pi(h))\ne f(\pi(h')).       \tag{M.9}
\]
The admissible histories and future queries in (M.9) must be specified.
Different auxiliary control modes carrying the same fact need not be
separated. This is a conditional design requirement, not a proved universal
decomposition or a learning theorem. A fixed finite register cannot retain
an unbounded number of independently queryable facts.

\subsection{What remains open, and what existing proofs need}
The difficult contribution is learning or synthesizing this decomposition
with an efficiently checkable future-error invariant and competitive cost.
A supplied decomposition plus a geometric error series is useful machinery,
but does not establish a general discovery theorem.

Theorem~\ref{III:M:audit:thm:addresses} and Proposition~\ref{III:M:audit:prop:checkpoint} do not
depend on solving the feedback, learned-congruence or general Hankel
decomposition targets. Their hypotheses and proofs are explicit above.
Lemma~\ref{III:M:audit:lem:guard} also does not assume that an inconclusive test can
always be resolved. Resolving those open targets is necessary only for the
corresponding stronger claims: general closed-loop agent memory, automatic
discovery of protected facts, or global optimality beyond the stated class.

\textbf{Proof obligations.} Specify the event/query language and reference
semantics; establish all-letter update compatibility; transport numerical and
model error to the declared answers; and prove the storage/update bound.
The ordinary counting-Hankel estimate does not automatically certify rare
conditional queries. The conditional bound divides by prefix mass.
A controlled agent whose actions change subsequent observations needs a
controlled-system or policy-level argument in addition to a fixed-input bound.

\textbf{Decisive experiment.} Use a frozen answering model and typed events
with a known oracle: insert facts, apply corrections and retractions, interleave
long irrelevant sequences, and ask delayed or counterfactual queries. Compare
an exact key--value baseline, retrieval/summary baselines, a homogeneous fading
state, the unequal-ratio state, and the persistent-plus-predictive construction.
Match total memory rather than just injected tokens. Report factual accuracy,
order sensitivity, certificate coverage/violations, update and readout latency,
and peak allocation. Afterwards test paraphrased natural-language interfaces,
separating extraction failures from memory failures. No such agent benchmark
is reported in this paper.

\section{Certified feedback, joint memory and consolidation}
\label{III:M:sec:agent-extensions}
This section gives conditional constructions for feedback, joint memory
and consolidation. All input alphabets, mode catalogues and structural constants
are fixed unless a bound explicitly varies them. The elementary packing,
coupling, quotient and concentration arguments below are not claimed as
new general principles. Their role is to state the exact executable contract.

\subsection{Cache decisions at the time of their certification}
Use the following time convention throughout this section:
\begin{equation}\label{III:M:eq:feedback-time}
 \sigma_j=\mathbf1\{g_{a_j}(x_{j-1})\ge0\},\qquad
 x_j=\tau_{a_j,r_{j-1}}(x_{j-1}),\qquad
 r_j=G(r_{j-1},a_j,\sigma_j).
\end{equation}
Both right-hand sides use the old mode; the guard uses the state before
the continuous update. A model whose guard uses the updated state requires
a different, explicitly stated convention. No guard is recomputed at an
intermediate replay state.

\begin{lemma}[Maximum stopped-window length]\label{III:M:lem:feedback-length}
For $0<r_a<1$, put $r_{\max}=\max_a r_a$ and
\[
 L_\theta=\left\lceil\frac{\log(1/\theta)}{\log(1/r_{\max})}\right\rceil.
\]
Every word stopped at its first product at most $\theta$, and every
internal startup word, has length at most $L_\theta$. For
$\theta=2^{-(p+c_0)}$, this is $O_H(p)$ for a fixed model.
\end{lemma}
\begin{proof}
The product of any $L_\theta$ ratios is at most
$r_{\max}^{L_\theta}\le\theta$; its first crossing occurs no later.
An internal word has product greater than $\theta$, so is shorter.
\end{proof}

The slowest contraction is decisive. For ratios $1/4,3/4$ and
$\theta=1/16$, the all-$3/4$ word first crosses at length $10$:
$(3/4)^9>1/16\ge(3/4)^{10}$. Replacing $r_{\max}$ by $r_{\min}$
would incorrectly give $2$.

\begin{theorem}[Certified branch-cache replay]\label{III:M:thm:branch-cache}
Assume the rational similarities, common compact $X$, initial point $x_0$,
and effective evaluators of Proposition~\ref{III:M:audit:prop:checkpoint}, but
replace its mode update by \eqref{III:M:eq:feedback-time}. Let $D\ge\operatorname{diam}(X)$.
Suppose at each incoming event a finite certified computation produces
an interval containing $g_{a_j}(x_{j-1})$ that excludes zero, using only
the deployed state, the new event and charged model data. Store its exact
sign bit $\sigma_j$. The hypothesis is an explicit restriction on the
admissible trajectories and guard algorithm, not an automatic consequence
of contraction.

Let $v_t=(a_t,\ldots,a_{t-\ell+1})$ be the stopped word and store
\[
 m_t=(b_t,I_t,\beta_t),\qquad
 b_t=r_{t-\ell},\quad I_t=\Mrank(v_t),\quad
 \beta_t=(\sigma_t,\ldots,\sigma_{t-\ell+1}).
\]
Replay the reversed letter--bit pairs from $(b_t,x_0)$ using the
simultaneous update $(\rho,y)\leftarrow
(G(\rho,a,\sigma),\tau_{a,\rho}(y))$.
Then the final mode is exactly $r_t$, and exact-arithmetic replay satisfies
\[
 \Mnorm{x_t-y_t}\le D r_{v_t}\le D\theta
\]
whenever deletion has occurred; if all history is retained, $y_t=x_t$.
Certified numerical replay adds its explicit $E_{\rm num}$.

On arrival of $a$, first certify the new bit on the current state enclosure,
then prepend $(a,\sigma)$, and trim letters and bits together. Advance $b_t$
through the dropped pairs in chronological order. This is an autonomous
update, including noninvertible mode resets. A fixed-width representation
uses at most
\begin{equation}\label{III:M:eq:branch-bits}
 \lceil\log_2|\mathcal R|\rceil+
 \lceil\log_2N_\theta\rceil+L_\theta
\end{equation}
persistent bits. The decoded address supplies the active bit-string length;
unused bits in the fixed field are zeroed. Addressing and mode replay keep
the previous polynomial bounds and $O_H(p)$ scratch. Total event time also
includes the guard's actual certification cost, for which no universal
bound is asserted here.
\end{theorem}
\begin{proof}
Initially the address and bit string are empty and $b_0=r_0$.
Inductively every cached bit is the reference event-time bit by the
certified enclosure. If $j$ oldest pairs are dropped after appending,
their forward mode transitions advance the base from $t-\ell$ to
$t-\ell+j=t+1-|v_{t+1}|$. No inverse transition is used.
Replaying cached pairs therefore follows exactly the reference mode path.
The true state at the retained base and $x_0$ both lie in $X$ and undergo
the same similarities. Their distance contracts by $r_{v_t}$. Before
deletion they start from the same $x_0$. Numerical errors follow the
geometric recurrence of the preceding checkpoint proposition.
Lemma~\ref{III:M:lem:feedback-length} bounds the extra bit field and the number
of added mode operations. Guard work is separate from these operations.
\end{proof}


The guard can use the image of $X$ under the cached mode path, enlarged
for numerical error. This outer enclosure need not be tight. At zero
margin the theorem does not promise a successful event: halt or abstention
must be declared, or a separate justified mechanism must supply the bit.
An external decision oracle is not an uncharged implementation.

\paragraph{Optional augmented-alphabet encoding.}
If only letters in a declared $\Sigma_{\rm gate}$ require a bit, encode
the alphabet consisting of one copy of each other letter and two copies
of each gated letter. Both copies have ratio $r_a$. The address counting
argument gives an upper rate $\widetilde s p+O_H(1)$, in addition to the
mode checkpoint, where
\[
 \sum_{a\notin\Sigma_{\rm gate}}r_a^{\widetilde s}
 +2\sum_{a\in\Sigma_{\rm gate}}r_a^{\widetilde s}=1.
\]
This is an upper bound on all encoded pair words. It is not an optimality
claim: many bit words can be unreachable and the copied geometric branches
need not be separated. Neither this construction nor \eqref{III:M:eq:branch-bits}
preserves the earlier coefficient $s$ without further information.

\subsection{Sound coalescence with its actual cardinality bound}
Here hypotheses are an outer representation of possible states, not a
probability distribution. Soundness means that the actual state remains
covered. It does not authorize selecting an arbitrary hypothesized mode
as an exact answer.

\begin{proposition}[A shared continuous address and a mode subset]
\label{III:M:prop:oneway-coalescence}
Suppose $\tau_{a,r}=\tau_a$ and the ratios satisfy the address assumptions.
Maintain the single address $I_t$ and a subset $S_t\subseteq\mathcal R$
containing the actual mode. For the new letter let $J_t\subseteq\{0,1\}$
contain every guard value on the sound continuous enclosure; if unresolved,
use $J_t=\{0,1\}$. Set
\[
 S_{t+1}=\{G(r,a,\sigma):r\in S_t,\ \sigma\in J_t\}.
\]
This representation remains sound and has at most $|\mathcal R|$ mode
labels at every time. A bit-mask implementation uses
$|\mathcal R|+\lceil\log_2N_\theta\rceil$ persistent bits.
\end{proposition}
\begin{proof}
The continuous path depends only on the input word, so every branch shares
its continuous address and sound error enclosure. The actual pair
$(r_t,\sigma_{t+1})$ occurs in the displayed Cartesian product; its successor
mode belongs to $S_{t+1}$. Taking a set merges duplicate labels and leaves
at most $|\mathcal R|$ elements. A common enclosure may discard correlations
and admit extra possibilities, but cannot discard the actual trajectory.
\end{proof}

The proposition controls representation size, not enclosure sharpness or
factual ambiguity. A singleton answer is certified only when all retained
possibilities justify it under the declared output contract.

\begin{proposition}[Separated gates: settled and transient counts]
\label{III:M:prop:dwell-coalescence}
Allow mode-dependent similarities with letter-only ratios. Let gated
events have binary successor choices; ordinary events keep the mode fixed.
Between consecutive gated events require at least $L_\theta$ ordinary
events. Start from the known initial state, or a settled representation
with at most one canonical stopped-window hypothesis per current mode.
Use common letter addressing and sound enclosures, and discard mode-path
details only after they leave the stopped window.

Then a sound implementation can maintain at most $|\mathcal R|$
hypotheses immediately before each gated event, and at most
$2|\mathcal R|$ at arbitrary intermediate times. Once the previous gate
has left the active window, histories with the same current mode can use
the common replay from $(r,x_0)$ and its $D\theta$ error enclosure.
This is canonical re-encoding with a certified error, not equality of
the original continuous states.
\end{proposition}
\begin{proof}
After $L_\theta$ ordinary events, the stopped word is contained wholly
in those events. Its mode is constant and equal to the current mode;
thus the mode at its base is the current mode. For each such mode, replay
from $x_0$ covers all true paths in that mode within $D\theta$, because
their base states lie in $X$. Common ratios make the retained letter word
identical across hypotheses. Coalescing to one canonical enclosure per
mode is therefore sound at this settled time.
The next gate creates at most two successors per settled hypothesis.
Ordinary events do not split them. The next gate occurs only after a new
settled time, which proves both counts by induction. One common address,
the remaining mode paths (at most one active gate each), and a finite
hypothesis catalogue suffice; their storage and guard work are charged.
\end{proof}


\paragraph{Why the all-time $|\mathcal R|$ bound is false.}
Take $\mathcal R=\{0,1\}$, $X=[0,1]$, ratio $1/4$ and
$\tau_r(x)=x/4+3r/4$. Before a gate, the enclosures
$[0,1/4]$ and $[3/4,1]$ are compatible with the two modes.
For $g(x)=(x-1/8)(x-7/8)$ both enclosures contain both signs.
With $G(r,a,\sigma)=\sigma$, splitting at the gate produces distinct
old-mode/new-mode pairs $(0,0),(0,1),(1,0),(1,1)$.
The new continuous enclosures from the two old modes lie respectively
in $[0,1/16]$ and $[15/16,1]$; sharing a new mode does not identify them.
There can therefore be four candidates for two modes immediately after
the gate, even though a sufficiently long subsequent dwell permits
canonical coalescence. This is a counterexample to the claimed inference
for the maintained sound hypothesis representation; it does not assert
that all its outer candidates must be reachable from one exact history.

\subsection{A matching lower bound for separated observable fibres}
\begin{theorem}[Conditional joint persistent-bit optimum]
\label{III:M:thm:joint-mode-optimum}
Use the letter-driven model and effective hypotheses of
Proposition~\ref{III:M:audit:prop:checkpoint}; this theorem does not include
unlogged continuous feedback. Let $\mathcal R_{\rm reach}$ be its finite
reachable mode set. Assume the closures $X_r$ of the reachable continuous
fibres are nonempty $s$-Ahlfors regular, with common positive regularity
constants and scales, where $\sum_a r_a^s=1$.
For actual reachable states assume
\[
 d_{\rm fut}((r,x),(r,y))\ge c\Mnorm{x-y},\qquad
 d_{\rm fut}((r,x),(r',y))\ge2\delta_R\quad(r\ne r'),
\]
with $c,\delta_R>0$. Assume also the uniform within-mode readout
$C$-Lipschitz bound and effective numerical error allocation of the
checkpoint proposition for every declared query. Futures are the
declared common input continuations. For sufficiently large $p$ with
$2^{-p}<\delta_R$,
\[
 \log_2|\mathcal R_{\rm reach}|+sp-O_H(1)
 \ \le B^*(2^{-p})\ \le
 \log_2|\mathcal R_{\rm reach}|+sp+O_H(1).
\]
Here $B^*$ minimizes deterministic autonomous persistent memory under
uniform future-answer TV error. Constants may depend on the fixed model,
including regularity, observability and readout constants.
\end{theorem}
\begin{proof}
Regularity gives, in each compact fibre closure, a packing at separation
$4\varepsilon/c$ of cardinality at least $a_H\varepsilon^{-s}$ at small
scales. Approximate its finitely many points by actual reachable states
within $\varepsilon/(2c)$; their separation remains greater than
$2\varepsilon/c$. Within a mode, observability makes the future distance
of any two selected histories greater than $2\varepsilon$. Across modes
the same conclusion follows from $\varepsilon<\delta_R$.
If two such histories shared a stored state, autonomous updates would
give identical stored states under every common continuation. The two
uniform approximation inequalities and TV's triangle inequality would
then give future distance at most $2\varepsilon$, a contradiction.
Consequently at least $|\mathcal R_{\rm reach}|a_H\varepsilon^{-s}$
codes are required. Taking logarithms proves the lower bound.

Every retained base mode is reachable. Indexing only those modes in the
checkpoint construction and choosing $\theta$ and numerical precision
for output error $\varepsilon$ gives the upper bound. Under a common
future input, equal modes remain equal and identical similarities do not
increase continuous discrepancies. The uniform query Lipschitz bound
therefore supplies the declared future contract as well as the current
readout contract. The original online addressing estimates apply.
\end{proof}

The upper Lipschitz hypothesis cannot be replaced by lower observability.
For example, the family of threshold queries $h_q(x)=\delta_{\mathbf1\{x\ge q\}}$
separates every distinct pair by TV $1$, however close their coordinates.
On an infinite reachable set no finite code achieves uniform error below
$1/2$ for that family. The theorem is a supplied-model result; it neither
discovers the fibres nor proves their regularity for a learned agent.
For a growing catalogue, dependence hidden in $O_H(1)$ must be made uniform.

\subsection{A conditional-query contract without a query prior}
\begin{proposition}[Uniform query kernels]\label{III:M:prop:query-kernels}
Let $h_q(r,x)$ be specified probability kernels for every declared query,
including queries outside the support of a sampling distribution. Suppose
\[
 \sup_q\MTV(h_q(r,x),h_q(r,y))\le C\Mnorm{x-y}
\]
for all admissible within-mode pairs. If the deployed mode is exact and
its state error is at most $e$, then its plug-in kernel has uniform query
error at most $Ce$, independently of any query prior.
\end{proposition}
\begin{proof}
Apply the assumed inequality to the exact and decoded states and take the
supremum over $q$. No conditioning of a joint word law is performed.
\end{proof}

For a null-probability query, the kernel is part of the specified reference
semantics, not an inferred conditional probability on a null event.
This proposition does not remove the rare-prefix obstruction for an
unconditional Hankel or completed-word TV approximation. It uses a
stronger, separately verified family of conditional readouts.

\subsection{External trajectory guarantees in a controlled environment}
Compare two agents with the same initial external-state distribution and
the same environment/observation kernels. The reference and deployed agents
receive the same history until their actions differ. Assume standard
measurable kernels (finite actions suffice below). ``Trajectory'' means
external states, observations and actions; it does not include the two
different internal memory representations. The memory error assumptions
below must hold on every shared history in the specified support, not just
on one realized test rollout.

\begin{theorem}[Guarded actions and probability of guard failure]
\label{III:M:thm:guarded-policy}
Let the reference policy choose the largest of finitely many scores
$g_a(r,x)$ with a fixed tie rule. At a shared history, suppose the mode is
exact and certified computed scores obey
$|\widetilde g_a-g_a(r,x)|\le\beta_t$ for every action, where for example
$\beta_t=L_g e_t+\eta_t$. Choose a computed maximizer $a^*$ only when
\[
 \widetilde g_{a^*}-\max_{b\ne a^*}\widetilde g_b>2\beta_t.
\]
Then the selected action equals the unique reference maximizer.
Let $\tau$ be the first failed action or memory certificate in a coupling
that keeps the two external trajectories identical before failure. If the
deployed policy is defined after failure as well, then
\[
 \MTV(P_{\rm ref}^{1:T},P_{\rm dep}^{1:T})
 \le\Pr(\tau\le T).
\]
In particular the laws coincide if certificates succeed almost surely
throughout the supported horizon. If failure at each next step, conditional
on no earlier failure, has probability at most $\alpha_t$, the right-hand
side is at most $\min(1,\sum_{t=1}^T\alpha_t)$.
\end{theorem}
\begin{proof}
For every $b\ne a^*$, the true score difference exceeds the computed
difference minus $2\beta_t$, hence is positive. Identical actions allow
the environment to use the same transition and observation samples; the
next external histories remain equal. Induction constructs the coupling
until $\tau$. The probability that its full external paths differ is at
most $\Pr(\tau\le T)$, which bounds TV. Summing the probabilities of
first failure gives the last assertion. Observed guard success on one
sample path does not bound this failure probability.
\end{proof}


\begin{theorem}[Stochastic-policy coupling]\label{III:M:thm:stochastic-policy}
Suppose on every shared history the two action kernels have TV distance
at most $d_t\le1$. Then
\[
 \MTV(P_{\rm ref}^{1:T},P_{\rm dep}^{1:T})
 \le1-\prod_{t=1}^T(1-d_t)
 \le\min\left(1,\sum_{t=1}^Td_t\right).
\]
In particular take $d_t=\min(1,L_\pi e_t)$ if exact modes and a uniform
$L_\pi$-Lipschitz stochastic policy are independently established on
shared histories. Uncertified mode errors are not covered by this choice.
\end{theorem}
\begin{proof}
While histories agree, maximally couple the next finite-action kernels;
their conditional probability of disagreement is at most $d_t$.
On agreement couple the environment identically. The probability of no
disagreement through time $T$ is at least $\prod_t(1-d_t)$. The coupling
inequality proves the first bound, and the union/product bound the second.
The environment need not be contractive: any amplification after a first
disagreement is already included in its probability.
\end{proof}

These theorems transfer a supplied uniform memory/policy certificate to
closed-loop behavior. They do not establish its premises for arbitrary
trained policies or infer total-variation guarantees from benchmark scores.

\subsection{Exact active-map quotients and certified approximate merging}
\begin{proposition}[Active-map quotient with a complete finite code]
\label{III:M:prop:active-map}
Let $K,V$ be fixed finite sets of sizes $k,v\ge1$. Events are deterministic
set/update, delete and no-op operations on $K\to V\cup\{\bot\}$.
Assume every future discrete output and transition depends on the past
only through the current map, and every continuous update depends only
on that map, its continuous state and the next event. No historical,
provenance, scheduled activation or expiration query is included unless
its required metadata is added to the state.
Then replacing the operation history by its active map preserves all
future behavior exactly. All maps can be encoded in
$\lceil\log_2 (v+1)^k\rceil$ bits. If every map is reachable and every
key including absence can be queried, this number is the exact minimum
fixed-width bit budget for the discrete map alone.

A simpler enumerative code for a map with $n$ active keys uses at most
\[
 \lceil\log_2(k+1)\rceil+
 \lceil\log_2\tbinom{k}{n}\rceil+n\lceil\log_2 v\rceil
\]
bits, including the variable cardinality. If $n$ is fixed and known as
part of the contract, the first term can be omitted.
\end{proposition}
\begin{proof}
Equal active maps and continuous states produce equal next maps and states
under every common event, and equal outputs. Induction proves exact
future equivalence. Each key has $v+1$ possibilities, yielding $(v+1)^k$
maps and a finite radix code. Any two different reachable maps differ at
a key; its exact query distinguishes them, proving the fixed-width lower
bound. Alternatively encode $n$, the ranked $n$-element key subset, and
its values in sorted-key order, which gives the displayed upper length.
\end{proof}

Canonical semantics does not imply that this serialization is shorter than
every previous representation. For $k=v=1$ there are two maps, empty and
present: omitting the cardinality from the conditional-size formula would
assign zero bits to both. An exact old/new codec selector costs one extra
bit and permits length $1+\min(L_{\rm old},L_{\rm map})$ if both decoders
and their data are available and charged. Conversion workspace, retained
logs and peak allocation must also be counted. Future-dated activations
cannot be erased merely because they are inactive now.

\begin{proposition}[Merging in a certified future pseudometric]
\label{III:M:prop:future-merge}
Let $d_{\rm fut}$ be the declared common-continuation pseudometric.
Suppose a possible true state $z$ belongs to at least one certified ball
$d_{\rm fut}(z,z_i)\le e_i$. Partition the candidate centres into groups
and choose admissible representatives $\bar z_j$ with certified distances
$d_{\rm fut}(z_i,\bar z_j)\le\delta_i$ for each member of group $j$.
Replacing the groups by their representatives with radii
$\max_{i\in j}(e_i+\delta_i)$ preserves sound coverage.

Under common deterministic inputs, if the transition is nonexpansive in
this future metric, successive certified replacements of radii at most
$\delta_1,\ldots,\delta_n$ increase a single-representative error by at
most $\sum_j\delta_j$. This sum, or a sharper proved transport bound,
must be charged against the total allowed error.
\end{proposition}
\begin{proof}
For the group containing a centre that covers $z$, the triangle inequality
gives $d_{\rm fut}(z,\bar z_j)\le e_i+\delta_i$. The representative thus
covers the true state with the stated radius. For repeated replacements,
apply this inequality after each common update; nonexpansiveness prevents
the intervening updates from enlarging the existing bound. In the usual
supremum-over-continuations metric, nonexpansiveness follows because
continuations beginning with a fixed letter form a subset of all
continuations. This does not compare agents taking different actions.
\end{proof}

The result certifies coverage, not a posterior weight or permission to
select an arbitrary remaining hypothesis as true. Computationally useful
certificates of the stated future distances remain a separate obligation.

\paragraph{When a midpoint estimate is justified.}
For an admissible convex within-mode domain and a proved bound
$d_{\rm fut}((r,x),(r,y))\le C\Mnorm{x-y}$ on all its points,
the midpoint of $x_i,x_j$ has future distance at most
$C\Mnorm{x_i-x_j}/2$ from either. This controls the centres only;
their existing uncertainty radii must be added. Convexity and the
\emph{future} bound are essential. A current-output Lipschitz constant
is insufficient when a later guard can switch a persistent mode.

\paragraph{A rational counterexample to Euclidean branch merging.}
Take $X=[-1,1]$, current mode $0$, and states $x=\pm\eta$, $0<\eta<1$.
At the next gated input set $r^+=\mathbf1\{x\ge0\}$ and $x^+=x/2$;
ordinary subsequent inputs preserve the mode. Querying that mode gives
opposite Dirac laws, hence future TV $1$. The continuous distance is
$2\eta\to0$. Merging into $x=0$ selects the positive branch and loses
the negative one. Even a perfectly contractive continuous map does not
make this merge safe. A nonconvex invariant set introduces the additional
problem that the midpoint need not be an admissible state.

\subsection{Verification from supplied probes}
Neither numerical closeness of trajectories nor an eigenvalue gap alone
establishes a statistical identification theorem. With exact noiseless
observations of $x_1=\lambda x_0$ and known $x_0=1$, a single observation
distinguishes $\lambda=1$ from every specified $\lambda<1$. Conversely,
hidden models with identical observable laws cannot be distinguished by
any sample size, even if their latent eigenvalues have a gap.

\begin{proposition}[Finite supplied probes under explicit sampling promises]
\label{III:M:prop:probe-test}
Fix $M$ supplied probes, $0<\alpha<1$, $0<\delta\le1$, $D\ge\delta$,
and $0<\gamma<1$. For each probe $i$ and chosen lag $L$, suppose one can
obtain $n$ independent identically distributed observations in $[0,1]$
with mean $\mu_i(L)$. Each probe is promised to be either persistent,
with $\mu_i(L)\ge\delta$, or fading, with
$0\le\mu_i(L)\le D(1-\gamma)^L$. Choose
\[
 L\ge\left\lceil\gamma^{-1}\log(4D/\delta)\right\rceil,
 \qquad
 n\ge\left\lceil8\delta^{-2}\log(2M/\alpha)\right\rceil.
\]
Classify a probe as persistent exactly when its sample mean exceeds
$\delta/2$. All $M$ classifications are correct with probability at
least $1-\alpha$. The count is $Mn$ probe observations at lag $L$,
not $n$ unrestricted training examples for an unknown memory architecture.
\end{proposition}
\begin{proof}
The lag makes every fading mean at most $\delta/4$, since
$(1-\gamma)^L\le e^{-\gamma L}$.
For a bounded observation $Y$ of mean $\mu$, convexity gives
$\mathbb E e^{tY}\le1-\mu+\mu e^t$. The logarithm of the centred
right-hand side has value and derivative zero at $0$ and second derivative
at most $1/4$ (the variance of a tilted Bernoulli variable). Hence
$\mathbb E e^{t(Y-\mu)}\le e^{t^2/8}$ for either sign of $t$.
Independence and exponential Markov, optimized at $t=4a$, give
$\Pr(|\bar Y-\mu|\ge a)\le2e^{-2na^2}$.
Taking $a=\delta/4$ and a union bound over $M$ probes gives failure
probability at most $2M e^{-n\delta^2/8}\le\alpha$.
On its complement the fading sample means are below $\delta/2$ and
persistent means above $3\delta/4$, so the decision rule is correct.
\end{proof}

The proposition assumes observable probes with separated means and
independent sampling. It does not identify arbitrary latent modes,
construct a protected fact congruence, validate a semantic extractor, or
prove persistence against every neutral word. Correlated trajectories
require their own concentration and effective-sample-size analysis.

\subsection{Nonconformality is not itself an impossibility theorem}
Take $A_0=A_1=\Mdiag(1/3,1/5)$ with translations $(0,0),(2/3,0)$.
The attractor is the middle-third Cantor set times $\{0\}$.
For $0\le s\le1$, every length-$n$ product has singular-value function
$3^{-ns}$. Both norm and singular-value pressures are therefore
$\log2-s\log3$, with the same zero $\log2/\log3<1$:
nonconformality does not imply a strict exponent gap.
With zero translations and distinct linear parts
$\Mdiag(1/3,1/5),\Mdiag(1/3,1/6)$, that pressure calculation remains valid
but the attractor is a singleton. Affinity dimension need not equal
covering dimension in unrestricted overlapping systems.
Memory rates require future-metric packing/covering bounds and effective
updates; dimension alone does not give $O_H(1)$ overhead. General
nonconformal constructions remain open here, not proved impossible.

\section{Mode-dependent ratios: pressure bounds and conditional optimality}\label{III:M:rev13:sec:mode}

Proposition ``Rational addresses with a mode checkpoint'' requires the ratio $r_a$ to depend
on the letter only. This section removes that restriction, gives a constructive pressure coefficient and a conditional matching
lower bound, and keeps the table-free polynomial update.

\subsection{Setting}
Let $\mathcal R$ be a finite mode set with deterministic transition
$G:\mathcal R\times\Sigma\to\mathcal R$ (resets allowed), initial mode $r_0$, and
\[
r_{t+1}=G(r_t,a_{t+1}),\qquad x_{t+1}=\tau_{a_{t+1},r_t}(x_t),\qquad
\tau_{a,m}(x)=r_{a,m}O_{a,m}x+t_{a,m},
\]
with rational ratios $r_{a,m}\in\mathbb Q\cap(0,1)$ depending on \emph{both} the letter
and the mode, orthogonal $O_{a,m}$, a common compact invariant set $X\ni x_0$ of known
diameter bound $D$, and all model data fixed. Write $r_{\min},r_{\max}$ for the extreme
ratios. Define the pressure matrix
\begin{equation}\label{III:M:rev13:eq:pressure}
M(s)_{m,m'}=\sum_{a\in\Sigma:\,G(m,a)=m'}r_{a,m}^{\,s},\qquad s\ge0 .
\end{equation}
Call the mode graph (edges $m\to G(m,a)$) \emph{strongly connected} if $M(0)$ is
irreducible. Then $s\mapsto\rho(M(s))$ is continuous and strictly decreasing, with
$\rho(M(0))=d\ge2$ and $\rho(M(s))\to0$, so there is a unique $s^\ast>0$ with
\begin{equation}\label{III:M:rev13:eq:sstar}
\rho(M(s^\ast))=1,\qquad M(s^\ast)u=u\ \text{ for some }u>0 .
\end{equation}
This is a symbolic pressure exponent. Its identification with geometric
Hausdorff dimension additionally requires a separation hypothesis, as in
\citet{mauldin_williams}. Overlapping or identical branches need not have
that dimension. In this section $d=|\Sigma|\ge2$ and $p\ge1$.

\subsection{The mode-aware stopped address}
For a history $a_1\cdots a_t$ with modes $r_0,\dots,r_t$, let
$\pi_j=\prod_{i=t-j+1}^{t}r_{a_i,r_{i-1}}$ be the product of the ratios \emph{actually
applied} by the $j$ newest letters ($\pi_0=1$). Let
$j^\ast=\min\{j:\pi_j\le\theta\}$ if this set is nonempty, else $j^\ast=t$. The stored state is
\[
m_t=(b_t,v_t),\qquad v_t=(a_t,\dots,a_{t-j^\ast+1})\ \text{(newest first)},\qquad
b_t=r_{t-j^\ast}.
\]

\begin{lemma}[Autonomous update and exact replay]\label{III:M:rev13:lem:mode-update}
From $(b_t,v_t)$ and a new letter $a$, replay $\operatorname{rev}(v_t)$ from $(b_t,x_0)$ to
obtain the modes $r_{t-j^\ast},\dots,r_t$ exactly, form $w=a\,v_t$, compute its
newest-first prefix products with these modes, cut at the first product $\le\theta$ (or
keep all of $w$ if none), and advance $b_t$ chronologically through the dropped oldest
letters. The result is $(b_{t+1},v_{t+1})$. Replaying $\operatorname{rev}(v_t)$ from
$(b_t,x_0)$ returns the exact current mode $r_t$ and a point $y_t$ with
\[
\Mnorm{x_t-y_t}\le D\,\pi_{j^\ast}\le D\theta\ \text{ if a deletion has occurred},\qquad
y_t=x_t\ \text{ otherwise}.
\]
Every retained word has length at most
$L_\theta=\lceil\log(1/\theta)/\log(1/r_{\max})\rceil$.
\end{lemma}
\begin{proof}
If $v_t$ is a leaf ($\pi_{j^\ast}\le\theta$) then $w$ has full product
$r_{a,r_t}\pi_{j^\ast}\le\theta$, so the first product $\le\theta$ among the newest-first
prefixes of the full new history occurs within $w$; these prefix products coincide with
those computed from the replayed modes, which are exact because $G$ is deterministic and
$b_t$ is the true mode at the window base. If $v_t$ is internal it equals the whole history,
so $w$ is the whole new history. Dropping the $j$ oldest letters moves the base from
$t-j^\ast$ to $t-j^\ast+j=t+1-|v_{t+1}|$ and $G$ is only evaluated forwards. The true state
at the window base and $x_0$ both lie in $X$ and undergo the same mode-dependent
similarities, whose composition has ratio $\pi_{j^\ast}$. The length bound is the
Lemma~\ref{III:M:lem:feedback-length}.
\end{proof}


\subsection{Counting the addresses}
For a mode $m$ and $x>0$ let
\[
W_m(x)=\#\Bigl\{c_1\cdots c_n\ (n\ge0)\text{ chronological from }m:\ 
\textstyle\prod_{i=1}^n r_{c_i,m_{i-1}}>x\Bigr\},
\]
where $m_0=m$ and $m_i=G(m_{i-1},c_i)$.
\begin{lemma}[Address set and growth]\label{III:M:rev13:lem:mode-count}
Identify a newest-first stored word with its chronological reversal for
counting and ranking. Every reachable stored state lies in
\[
\begin{aligned}
\mathcal A_\theta={}&\{(r_0,w):\ w\text{ internal from }r_0\}\\
&\cup\bigcup_{b\in\mathcal R,\ c\in\Sigma}\{(b,\,c\,w'):\ w'\text{ chronological from }
G(b,c)\text{ with product in }(\theta,\theta/r_{c,b}]\},
\end{aligned}
\]
so that $|\mathcal A_\theta|=W_{r_0}(\theta)+\sum_{b,c}\bigl[W_{G(b,c)}(\theta)-
W_{G(b,c)}(\theta/r_{c,b})\bigr]$. If the mode graph is strongly connected then, with
$C_u=\max_mu_m/\min_mu_m$ for the Perron vector in \eqref{III:M:rev13:eq:sstar},
\[
W_m(x)\le\frac{C_u\,r_{\min}^{-s^\ast}}{d-1}\,x^{-s^\ast}\quad(m\in\mathcal R,\ 0<x<1),
\qquad
|\mathcal A_\theta|\le\frac{(1+d|\mathcal R|)\,C_u\,r_{\min}^{-s^\ast}}{d-1}\,\theta^{-s^\ast}.
\]
\end{lemma}
\begin{proof}
A leaf $v_t=(c_{j^\ast},\dots,c_1)$ with base $b$ and oldest letter $c_1$ satisfies
$\pi_{j^\ast}\le\theta<\pi_{j^\ast-1}$, and $\pi_{j^\ast}=r_{c_1,b}\,\pi'$ where $\pi'$
is the product of $c_2\cdots c_{j^\ast}$ read from $G(b,c_1)$; thus
$\pi'\in(\theta,\theta/r_{c_1,b}]$. Shorter newest suffixes have larger products, so no
further condition is needed. Internal states retain the whole history, so their base is
$r_0$. This proves the inclusion and the count.

For the growth, consider the forward stopped tree from $m$ at threshold $x$ (a branch
stops at its first node of product $\le x$). It is a full $d$-ary tree whose internal
nodes are counted by $W_m(x)$, so it has $(d-1)W_m(x)+1$ leaves, each of product in
$(r_{\min}x,x]$. Induction from the leaves upward, using $M(s^\ast)u=u$ at every internal
node of product $R$ and mode $m'$ (its children contribute
$\sum_aR^{s^\ast}r_{a,m'}^{s^\ast}u_{G(m',a)}=R^{s^\ast}u_{m'}$), gives
$\sum_{\text{leaves }\lambda}\rho_\lambda^{s^\ast}u_{\mathrm{end}(\lambda)}=u_m$.
Hence $(\#\text{leaves})(r_{\min}x)^{s^\ast}\min_mu_m\le\max_mu_m$, which gives the bound
on $W_m(x)$; the bound on $|\mathcal A_\theta|$ follows from the count with
$W_{G(b,c)}(\theta/r_{c,b})\ge0$.
\end{proof}


\begin{lemma}[Polynomial counting without listing words]\label{III:M:rev13:lem:mode-dp}
For a fixed model, every $W_m(x)$ with $x=\theta/R$, $R$ a product of ratios, and every
rank/unrank operation on $\mathcal A_\theta$ can be computed with $O_H(p^{d|\mathcal R|+O(1)})$
bit operations on $O_H(p)$-bit integers, using exact rational comparisons only.
\end{lemma}
\begin{proof}
The product of a chronological word depends only on its type vector
$n\in\MN^{\Sigma\times\mathcal R}$, $n_{a,m'}$ being the number of positions at which
letter $a$ is read in mode $m'$: the product is $\prod r_{a,m'}^{n_{a,m'}}$. Let
$N(n,m,m'')$ be the number of words from $m$ of type $n$ ending in mode $m''$. Then
$N(0,m,m'')=[m=m'']$ and
\[
N(n,m,m'')=\sum_{(m',a):\,G(m',a)=m'',\ n_{a,m'}\ge1}N(n-e_{a,m'},m,m').
\]
Words have length at most $L_\theta=O_H(p)$, so there are
$\binom{L_\theta+d|\mathcal R|}{d|\mathcal R|}=O_H(p^{d|\mathcal R|})$ type vectors, each
with a count of $O_H(p)$ bits, and
$W_m(x)=\sum_{n:\ \prod r^{n}>x}\sum_{m''}N(n,m,m'')$. Rank and unrank use preorder in the
forward tree from $G(b,c)$: the subtree below a node of product $R$ and mode $m'$ has
$W_{m'}(\theta/R)$ internal nodes, of which $W_{m'}(\theta/R)-W_{m'}(\theta/(r_{c,b}R))$
have product in the required window; summing the sizes of earlier sibling subtrees along
the $O_H(p)$ letters of the address gives the ordinal, and the inverse operation chooses
children by comparison. Offsets for $(b,c)$ blocks are the corresponding block counts.
\end{proof}


\begin{proposition}[Streaming type counts with linear scratch]\label{III:M:rev13:prop:euler-count}
Let $n=(n_{a,m})$ have total length $L$, and let $N(n,s,t)$ count
chronological words from $s$ to $t$ with those edge multiplicities.
Let $d_m=\sum_a n_{a,m}$ and $d_m^{\rm in}=\sum_{(a,b):G(b,a)=m}n_{a,b}$.
For $L>0$ first require
$d_m-d_m^{\rm in}=[m=s]-[m=t]$. Restrict to the active vertices
(incident to a positive edge, together with $s,t$), and let $T_t(n)$ be
the weighted number of directed spanning trees oriented toward $t$,
where an edge $(m,a)$ has weight $n_{a,m}$ and loops cannot enter a tree.
If the balance condition fails, or an active $m\ne t$ has $d_m=0$,
then $N=0$. Otherwise
\[
 N(n,s,t)=T_t(n)\,
 \frac{\displaystyle\prod_m\binom{d_m}{(n_{a,m})_{a\in\Sigma}}}
      {\displaystyle\prod_{m\ne t}d_m}.
\]
The empty-word value is $N(0,s,t)=[s=t]$.
For fixed model size this formula can be evaluated in polynomial bit time
using $O_H(p)$ scratch when $L=O_H(p)$. Consequently streaming the type
vectors gives polynomial-time rank/unrank and updates with linear scratch,
without memoising a polynomial-size DP table.
\end{proposition}
\begin{proof}
Label temporarily all copies of every edge. In an Euler trail ending at $t$,
the last exit from each vertex other than $t$ forms an arborescence toward
$t$. Conversely, choose such a tree, place its edge last in the outgoing
list at each non-root vertex, and order all remaining outgoing edges.
Following these lists from $s$ uses every edge: if it stopped too early,
the balance equations force a stop at $t$. The remaining edge degrees
are then balanced. Any unused outgoing list has its tree edge unused;
following these last-exit edges, the remaining degree balances would force
an unused exit at $t$, contradicting the stop. Thus the labelled count is
$T_t(n)\prod_m(d_m-1+[m=t])!$.
Each unlabelled word has exactly $\prod_{a,m}n_{a,m}!$ labelled versions.
Dividing and rewriting factorial ratios as multinomials proves the formula,
including the case $s=t$. A zero tree count gives zero walks.

For fixed $|\mathcal R|$, $T_t(n)$ can be evaluated by explicitly summing
all finitely many choices of one outgoing edge type at every non-root
vertex and testing whether they form a tree. Each summand is a product of
at most $|\mathcal R|-1$ integers bounded by $L$. This avoids any reliance
on a numerical determinant. Multinomials are formed as products of
binomial coefficients using exact incremental divisions. Their product is
at most $d^L$; the tree factor and denominator have only $O_H(\log L)$
bits. Hence every retained integer has $O_H(L)$ bits. In contrast, storing
uncancelled $L!$ would require $\Theta(L\log L)$ bits.
Enumerating the $O_H(p^{d|\mathcal R|})$ type vectors needs only the
current vector, count and sum. The sums $W_m(x)$ and the subtree differences
in Lemma~\ref{III:M:rev13:lem:mode-dp} therefore use $O_H(p)$ scratch. Recomputing them
along an $O_H(p)$-letter path preserves polynomial time. With elementary
integer arithmetic an ample bound is $O_H(p^{d|\mathcal R|+6})$ per
rank/unrank or update. This exponent depends on the fixed model and is not
a uniform low-degree bound.
\end{proof}



\subsection{The theorem}
\begin{theorem}[Pressure upper bound and conditional persistent-bit optimum]\label{III:M:rev13:thm:mode}
Assume the setting above, a strongly connected mode graph, and the effective evaluation
hypotheses of Proposition ``Rational addresses with a mode checkpoint'' (certified
$q$-bit evaluators for $x_0$ and the mode-dependent affine data, a readout uniformly
$C$-Lipschitz in the continuous coordinate within each mode). Let $s^\ast$ be defined by
\eqref{III:M:rev13:eq:sstar} and $\theta=2^{-(p+c_0)}$ with $CD2^{-c_0}\le1/2$.
\begin{enumerate}[label=(\roman*)]
\item The memory $(b_t,\operatorname{rank}(v_t))$ of Lemma~\ref{III:M:rev13:lem:mode-update} is an
autonomous finite-state executable with exact mode recovery, output error $2^{-p}$, and
\[
B(p)\le s^\ast p+\lceil\log_2|\mathcal R|\rceil+O_H(1)
\]
persistent bits; its update and readout are polynomial in $p$ (Lemma~\ref{III:M:rev13:lem:mode-dp}
plus the stipulated evaluators) with $O_H(p)$ scratch by Proposition~\ref{III:M:rev13:prop:euler-count} and streaming replay. No exponential
word or transition table is stored.
\item Let $\mathcal R_{\rm reach}$ be the reachable modes and assume, as in Theorem
``Conditional joint persistent-bit optimum'', that the closures $X_r$ of the reachable
continuous fibres are $s^\ast$-Ahlfors regular with common constants, that
$d_{\rm fut}((r,x),(r,y))\ge c\Mnorm{x-y}$ within modes and
$d_{\rm fut}((r,x),(r',y))\ge2\delta_{\mathcal R}$ across modes. Then for
$2^{-p}<\delta_{\mathcal R}$,
\[
\log_2|\mathcal R_{\rm reach}|+s^\ast p-O_H(1)\ \le\ B^\ast(2^{-p})\ \le\
\log_2|\mathcal R_{\rm reach}|+s^\ast p+O_H(1).
\]
\item $s^\ast\le\bar s$, where $\bar s$ solves $\sum_a\bar r_a^{\,\bar s}=1$ with
$\bar r_a=\max_mr_{a,m}$ a pointwise dominating letter-only ratio; the inequality is strict whenever some ratio depends on the mode along a cycle of
the mode graph.
\end{enumerate}
\end{theorem}
\begin{proof}
(i) Lemma~\ref{III:M:rev13:lem:mode-update} gives exact modes and continuous error $\le D\theta$,
hence output error $\le CD\theta\le2^{-p-1}$ by the choice of $c_0$; the certified
replay error allocation of Proposition~\ref{III:M:audit:prop:checkpoint} contributes the other
$2^{-p-1}$ because only the mode-dependent constants change. The address ordinal ranges over
$\mathcal A_\theta$, so Lemma~\ref{III:M:rev13:lem:mode-count} gives
$\lceil\log_2|\mathcal A_\theta|\rceil\le s^\ast(p+c_0)+\log_2\frac{(1+d|\mathcal R|)C_u}{d-1}
+s^\ast\log_2\frac1{r_{\min}}+1$; the base mode is stored separately or folded into the
ordinal. Lemma~\ref{III:M:rev13:lem:mode-dp} gives the costs.
(ii) Indexing only reachable base modes in (i) gives the upper bound. For the lower bound
the proof of Theorem ``Conditional joint persistent-bit optimum'' applies verbatim: in
each reachable fibre closure, Ahlfors regularity of exponent $s^\ast$ yields a
$(4\varepsilon/c)$-separated set of cardinality $\ge a_H\varepsilon^{-s^\ast}$, approximated
by reachable states of the same mode within $\varepsilon/(2c)$; within-mode observability
and the inter-mode margin make all these histories pairwise future-separated by more than
$2\varepsilon$, and behavioural packing forces $|\mathcal R_{\rm reach}|a_H\varepsilon^{-s^\ast}$
distinct codes.
(iii) $M(s)\le\overline M(s)$ entrywise, where
$\overline M(s)_{m,m'}=\sum_{a:G(m,a)=m'}\bar r_a^{\,s}$ has all row sums
$\sum_a\bar r_a^{\,s}$, equal to $1$ at $s=\bar s$; hence
$\rho(M(\bar s))\le\rho(\overline M(\bar s))=1$ and $s^\ast\le\bar s$. If some
$r_{a,m}<\bar r_a$ on an edge of a cycle, irreducibility makes the inequality strict.
\end{proof}


\begin{remark}[Separation and startup are substantive hypotheses]
The invariant graph-directed fibres satisfy
\[
 K_r=\bigcup_{(m,a):G(m,a)=r}\tau_{a,m}(K_m).
\]
Their construction matrix is $M(s)^\top$, which has the same spectral
radius. A mass construction for these incoming cylinders uses a positive
Perron vector of $M(s^\ast)^\top$, rather than the outgoing-tree vector
$u$ used in Lemma~\ref{III:M:rev13:lem:mode-count}. With incoming strong separation
and strong connectivity the usual cylinder argument gives
$s^\ast$-regular fibres; compare \citet{mauldin_williams}.
If $x_0\in K_{r_0}$, these fibres are the closures of reachable states:
to approximate a chosen length-$n$ incoming cylinder, first reach its base
mode from $r_0$ and then follow that cylinder; contraction makes the initial
point immaterial as $n\to\infty$.

For $x_0\notin K_{r_0}$ the extra startup orbit need not be finite. For
example the one-mode Fibonacci system at $x_0=3/10$ has infinitely many
reachable points $(3/10)4^{-n}$ outside $K$. They belong to scaled copies
of the gap $(1/4,1/2)$. One must therefore not infer Ahlfors regularity of
the entire reachable closure from ``finitely many startup points''.
For the lower bound, it also suffices to assume regular observable compact
subsets of the reachable fibre closures, rather than regularity of the
whole closures.

Without separation/observability there is no matching lower bound from
pressure alone. For one mode with two identical maps $\tau_a(x)=x/2$ and
$x_0=0$, $s^\ast=1$ but one state reproduces the continuous state exactly.
\end{remark}


\begin{remark}[Reducible mode graphs]
Without strong connectivity let $s^\ast=\sup\{s:\rho(M(s))\ge1\}$ over the modes reachable
from $r_0$. For every $s>s^\ast$ the vector $u=(I-M(s))^{-1}\mathbf 1>0$ satisfies
$M(s)u\le u$, and the induction of Lemma~\ref{III:M:rev13:lem:mode-count} gives
$W_m(x)\le C_s x^{-s}$, hence $B(p)\le sp+O_{H,s}(1)$ for every $s>s^\ast$; two
components with the same exponent in sequence can add a $\log p$ term, so the clean
$O_H(1)$ statement is the irreducible one.
\end{remark}



\begin{remark}[What this does not settle]
Feedback-driven modes (guards on the continuous state) still need the branch cache or a
certified reconstruction; the mode-dependent counting above assumes letter-driven $G$,
exactly as the checkpoint proposition. Determining the optimal branch-cache coefficient in the general feedback
setting remains open.
\end{remark}

\section{From exact algebra to a finite executable}\label{III:U:sec:execution}
The general certificate returns a realization in exact algebraic arithmetic.
Separate the three error contributions as follows:
\[
 \sup_w|f(w)-\widehat g(w)|
 \le\epsilon_{\rm reduction}+\epsilon_{\rm coefficients}
                                      +\epsilon_{\rm execution}.
\]
The two linear errors can be certified using the joint Stein construction.
The rounded executable has a separate nonlinear state recurrence.

\subsection{A sufficient integer-execution certificate}
Use a row state initialized at $c^*$, updated by right multiplication by
$T_a$, and read against $b$; this respects the product order in
$g(w)=c^*T_wb$. The following standard positive-majorant bound is a useful
reference certificate.
\begin{proposition}[Componentwise execution tube]\label{III:U:prop:tube}
Let $M_{ij}=\max_a|(T_a)_{ij}|$ and suppose $\rho(M)<1$.
Use column vectors for componentwise absolute values. If each machine
step incurs local state error at most $\eta\ge0$, and the initial error
is at most $e_0$, then
\[
 e_t\le (M^\top)^te_0+\sum_{j=0}^{t-1}(M^\top)^j\eta.
\]
If $e_0\le E=(I-M^\top)^{-1}\eta$, then $e_t\le E$ for all $t$,
and the state contribution to output error is at most $|b|^\top E$.
An invariant state-radius vector is
$R=(I-M^\top)^{-1}(|x_0|+\eta)$.
\end{proposition}
\begin{proof}
Taking componentwise absolute values after one update gives
$e_{t+1}\le M^\top e_t+\eta$. Iterate and use the nonnegative Neumann
series for the inverse. Also $R\ge|x_0|$ and
$M^\top R+\eta=R-|x_0|\le R$, proving the state invariant.
The output bound is the triangle inequality.
\end{proof}

Unbounded verification of quantized filters has established precedents
\citep{cox2014}. The local bound must cover the whole invariant region, including products,
accumulation, coefficient representation, and output evaluation. For a fixed
integer kernel, the accumulator width is checked against all products and
sums on that region. Tie handling and operation order are part of the
machine specification. Saturation is a different machine and requires its
own proof. Initial state and readout quantization are charged explicitly.

This sufficient condition is stronger than mean-square stability and may
fail for a stable source. It is not a general replacement for the Hankel
certificate. Nor is its geometric error accumulation a new principle.
Searching for a useful invariant, choosing a base, and allocating bits are
separate optimization tasks.

\section{Polynomial transfer computation: the remaining completeness gap}\label{III:U:sec:oracle}
Positive contracting two-state predictive models represented by homographic
maps admit a different computational question: when can word enumeration in
a causal state encoder be replaced by a polynomial transfer calculation?
The following theorem gives a sufficient analytic certificate. Its existence
does not follow from the exact-bit family.

\subsection{A two-radius analytic certificate}
On normalized coordinates $x\in[-1,1]$, write
\[
 G_b(x)=\frac{\alpha_b+\beta_bx}{1+\xi_bx},\qquad
 g_b(x)=\kappa_b(1+\xi_bx)^{-2},\qquad 0<\kappa_b<1.
\]
Suppose rational $1<w<R$ satisfy, for each branch,
\[
 |\xi_b|R<1,\qquad \max(|G_b(-R)|,|G_b(R)|)\le w.
\]
A real homography maps the boundary circle to a circle centered on the
real axis; its real endpoints give the maximum modulus. Thus
$G_b(\overline D_R)\subset\overline D_w$. Strict invariance on the
real interval and poles outside it guarantee that such radii exist near
one, though their quantitative margins remain model-dependent.

For $\sigma\in[s_-,s_+]\subset(0,\infty)$ define
\[
 (\mathcal L_\sigma P)(x)=\sum_{b=0}^1g_b(x)^\sigma P(G_b(x)).
\]
Use the analytic logarithm equal to zero at $x=0$ for $1+\xi_bx$.
With $A_* =\lceil2s_+\rceil$ one may take
\[
 W=\left\lceil\sum_b(1-|\xi_b|R)^{-A_*}\right\rceil,
 \qquad\Lambda=\frac W{1-w/R},\qquad
 \delta_m=\frac{WR^{-m}}{R-1}.
\]
For $\|P\|_w=\sum_j|c_j|w^j$ and coefficient truncation $\Pi_m$,
Cauchy's bound gives
\[
 \|\Pi_m\mathcal L_\sigma P\|_w\le\Lambda\|P\|_w,
 \qquad\|(I-\Pi_m)\mathcal L_\sigma P\|_{\infty,[-1,1]}
                                      \le\delta_m\|P\|_w.
\]
The section matrix can be formed without enumerating words. Its entries are
\[
 M_{kj}=\sum_b\kappa_b^\sigma
 \sum_{h=0}^{\min(j,k)}\binom jh\alpha_b^{j-h}\beta_b^h
 (-\xi_b)^{k-h}\frac{(2\sigma+j)_{k-h}}{(k-h)!}.
\]
For one branch, compute column zero recursively and subsequent columns by
\[
 a_{k+1}^{(0)}=-\xi_b\frac{2\sigma+k}{k+1}a_k^{(0)},\qquad
 a_k^{(j+1)}=\alpha_ba_k^{(j)}+\beta_ba_{k-1}^{(j)}
                                      -\xi_ba_{k-1}^{(j+1)}.
\]
This costs $O(m^2)$ scalar operations. Real powers and rounding still
require certified enclosures in a finite implementation.

\begin{theorem}[Correct propagation of repeated truncations]\label{III:U:thm:oracle}
Assume the real branch maps satisfy $G_b([-1,1])\subseteq[-1,1]$ for $b=0,1$, and the two-radius analytic bounds above hold for fixed $1<w<R$.
Let $V_n=\mathcal L_\sigma^n1$, $P_0=1$, and
$P_{n+1}=\Pi_m\mathcal L_\sigma P_n$. Assume also $g_b(x)\le1$
on $[-1,1]$. With $K=\max(2,\Lambda)$,
\[
 \|V_n-P_n\|_\infty\le n\delta_mK^{n-1}.
\]
If $g_b(x)\ge g_->0$, $g_-<1$, put $a_*=2g_-^{s_+}$.
Then $V_n\ge a_*^n$, and relative error at most $\eta$ follows when
\[
 R^m\ge\frac{nWK^{n-1}}{(R-1)\eta a_*^n}.
\]
Consequently $m=O_H(n+\log(1/\eta)+\log(n+1))$ suffices for fixed
model $H$ and its certified margins.
\end{theorem}
\begin{proof}
The polynomial norm grows at most as $\Lambda^n$.
Positivity and the two branches give
$\|\mathcal L_\sigma f\|_\infty\le2\|f\|_\infty$.
Hence $E_{n+1}\le2E_n+\delta_m\Lambda^n$, with $E_0=0$.
Summing
$\sum_{j=0}^{n-1}2^{n-1-j}\Lambda^j\le n\max(2,\Lambda)^{n-1}$
yields the displayed bound. Each real branch weight is at least
$g_-^{s_+}$, so induction gives $V_n\ge a_*^n$.
\end{proof}


The distinction $M^n1\ne\Pi_m\mathcal L_\sigma^n1$ is essential:
truncation and iteration generally do not commute. The propagation bound
above is therefore necessary.

\subsection{What is established and what remains conditional}
The exact-arithmetic theorem removes exponential word enumeration for this
finite-horizon transfer calculation. In the intended regime
$n=O_H(\log p)$, it gives degree $O_H(\log p)$ under the stated relative
accuracy requirements. It does not by itself prove that a numerical
candidate generator always returns a polynomial accepted by the downstream
rational encoder certificate.

The desired $O_H(p^2\log p)$ bit cost depends on fast arithmetic and
specific subroutine cost bounds. Ordinary arithmetic gives a different
bound. The compiler, stored polynomial, precision, and online update
must all be charged. The unresolved obligation is a terminating, uniformly certified construction over the
declared model class with quantitative constants, including a proof that
the proposed acceptance conditions can always be met. Successful instances
do not establish that obligation.

\paragraph{A second logical boundary: dimension versus bounded overhead.}
A box dimension $h$ generally gives
$\log N(\epsilon)=h\log(1/\epsilon)+o(\log(1/\epsilon))$, not an
$O(1)$ remainder. For $K=\{0\}\cup\{2^{-j}:j\ge1\}$, the dimension
is zero but the covering number at scale $2^{-p}$ grows on the order of
$p$, giving $\log_2N\asymp\log_2p$. Stronger two-sided covering or
separation hypotheses are required for a constant-overhead memory law.
\subsection{Certified arithmetic error without a completeness claim}
\label{III:U:sec:transfer-rounding}
Strictly contracting branches do not make an unnormalized transfer
operator a contraction toward a positive fixed function.
For example, take
\[
 G_0(x)=x/4-1/2,\quad G_1(x)=x/4+1/2,\quad
 g_0=g_1=1/4,\quad \sigma=1.
\]
These branches are separated on $[-1,1]$ and satisfy the two-radius
assumptions with $R=2$, $w=3/2$, and $\xi_b=0$. Nevertheless,
$\mathcal L_\sigma1=1/2$ and $V_n=2^{-n}$, so there is no positive
limit of $V_n$. A positive
eigenfunction is not, without normalization, a fixed function.
Even after normalizing the leading eigenvalue, a quantitative
convergence rate and its normalization must be proved separately.

\begin{theorem}[Finite-horizon rounding and conditional eigenfunction error]
\label{III:U:thm:transfer-rounding}
Under Theorem~\ref{III:U:thm:oracle}, let
$A_m=\Pi_m\mathcal L_\sigma$ and compute degree-at-most-$m$
polynomials by
\[
 \widehat P_0=1,\qquad
 \widehat P_{j+1}=A_m\widehat P_j+e_j,\qquad
 \|e_j\|_w\le r_j.
\]
The certified errors $e_j$ must include matrix-coefficient approximation,
arithmetic rounding, and serialization. Then
\begin{equation}
 \|V_n-\widehat P_n\|_{\infty,[-1,1]}
 \le n\delta_m K^{n-1}
                   +\sum_{j=0}^{n-1}\Lambda^{n-1-j}r_j.
 \label{III:U:eq:transfer-rounded-error}
\end{equation}
If, in addition, a specified positive function $h\ge h_{\min}>0$
and constants $C<\infty$, $0<\vartheta<1$ satisfy the independently
proved normalization and convergence bound
$\sup_x|V_n(x)-h(x)|/h(x)\le C\vartheta^n$, then
\begin{equation}
 \sup_x\frac{|\widehat P_n(x)-h(x)|}{h(x)}
 \le C\vartheta^n+
 \frac{n\delta_mK^{n-1}
                 +\sum_{j=0}^{n-1}\Lambda^{n-1-j}r_j}{h_{\min}}.
 \label{III:U:eq:transfer-conditional-eigenfunction}
\end{equation}
No such convergence assumption is inferred from branch contraction alone.
\end{theorem}
\begin{proof}
For the exact truncated iteration $P_j=A_m^j1$, repeated substitution
gives $\widehat P_n-P_n=\sum_{j=0}^{n-1}A_m^{n-1-j}e_j$.
Since $\|A_m\|_{w\to w}\le\Lambda$ and
$\|P\|_{\infty,[-1,1]}\le\|P\|_w$, its norm is bounded by the
sum in \eqref{III:U:eq:transfer-rounded-error}. Add
Theorem~\ref{III:U:thm:oracle}'s finite-horizon error. Dividing by
$h_{\min}$ and adding the separately assumed horizon-convergence error
proves \eqref{III:U:eq:transfer-conditional-eigenfunction}.
\end{proof}


An $O(m^2)$ construction of the section matrix does not account for
$n$ repeated dense matrix--vector products, which cost $O(nm^2)$
scalar operations without further structure. Bit precision, evaluation
of real powers, interval dependency, and verified inversion also need
their own bounds. A small polynomial approximation error is not by
itself proof of acceptance by the complete state compiler.
Accordingly, universal compiler termination, a sufficient precision rule, and uniform offline/online bit bounds remain unestablished.


\subsection{Normalized transfer convergence on a strict log-Lipschitz cone}
\label{III:U:sec:normalized-transfer}
The endpoint value $L=c/(1-\gamma)$ gives
only cone invariance. Strict inclusion requires a strictly larger
$L$. Moreover the relevant order is that of the regularity cone,
not only pointwise positivity in $C([-1,1])$.

\begin{theorem}[Quantitative normalized positive eigenfunction]
\label{III:U:thm:normalized-transfer}
On $I=[-1,1]$, let
$(\mathcal L f)(x)=\sum_{a=1}^d w_a(x)f(G_a(x))$, where
$G_a:I\to I$ have Lipschitz constants at most $\gamma<1$,
$0<w_-\le w_a(x)\le w_+$, and
$|\log w_a(x)-\log w_a(y)|\le c|x-y|$.
Choose $L>c/(1-\gamma)$ and put $L'=c+\gamma L<L$.
Define
\[
 t_* =\frac{L-L'}{L+L'}e^{-2(L+2L')},\qquad
 \Delta=2\log(1/t_*),\qquad \vartheta=\tanh(\Delta/4)<1.
\]
There is a unique eigenfunction $h>0$ in the log-Lipschitz cone of
constant $L$, normalized by $h(0)=1$. It satisfies
\[
 \mathcal L h=\lambda h,\quad
 e^{-L'}\le h(x)\le e^{L'},\quad
 d w_-e^{-L'}\le\lambda\le d w_+e^{L'}.
\]
Writing $V_n=\mathcal L^n1$, for $n\ge1$,
\begin{equation}
 \sup_{x\in I}\left|\frac{V_n(x)/V_n(0)}{h(x)}-1\right|
       \le e^{\Delta\vartheta^{n-1}}-1.
 \label{III:U:eq:normalized-transfer}
\end{equation}
\end{theorem}
\begin{proof}
Let $K_L$ be the closed convex cone of nonnegative Lipschitz functions
with $f(x)\le e^{L|x-y|}f(y)$ for all $x,y$; its only function
vanishing at a point is zero. Each positive summand of $\mathcal L f$
has log-Lipschitz constant at most $c+\gamma L$, and their sum has
the same bound. Thus $\mathcal L(K_L\setminus\{0\})\subset K_{L'}$.

For $f,g\in K_{L'}$ normalized at zero, $g(y)/f(y)\ge e^{-2L'}$,
and hence also $g(y)/f(y)\ge e^{-4L'}$. For $r=|x-y|\in(0,2]$,
\[
 \frac{e^{Lr}-e^{L'r}}{e^{Lr}-e^{-L'r}}
 \ge\frac{L-L'}{L+L'}e^{-2L}.
\]
This follows by expressing numerator and denominator as integrals
of $r e^{ur}$. Therefore
$e^{Lr}g(y)-g(x)\ge
t_*[e^{Lr}f(y)-f(x)]$. Pointwise positivity also holds, so
$g-t_*f\in K_L$, and symmetrically $f-t_*g\in K_L$.
The image diameter in the Hilbert metric of $K_L$ is at most
$2\log(1/t_*)=\Delta$.

For projective contraction, use the Lipschitz Banach space: the strict
constant $L'<L$ places these positive image functions in the interior
of $K_L$. The normalized slice $f(0)=1$ of $K_L$ is convex and compact in the
uniform norm by uniform bounds and equicontinuity. The continuous
map $f\mapsto\mathcal L f/(\mathcal L f)(0)$ has a fixed point
by Schauder. Birkhoff's projective contraction theorem
\citep{gaubert2015} gives contraction factor at most $\vartheta$ on
comparable rays. Every eigenfunction in $K_L$ lies in the image,
whose diameter is finite, so this also gives uniqueness.
Both $\mathcal L1$ and $\mathcal Lh$ are image points; iteration
gives projective distance at most $\Delta\vartheta^{n-1}$.
Cone order implies pointwise order. Since both representatives have
value one at zero, their pointwise ratios are between $e^{-D}$ and
$e^D$ when their projective distance is at most $D$. This proves
\eqref{III:U:eq:normalized-transfer}. The fixed point lies in $K_{L'}$;
its normalization gives the bounds on $h$, and evaluation at zero
gives those on $\lambda$.
\end{proof}


\begin{corollary}[A certified polynomial eigenfunction approximator]
\label{III:U:cor:normalized-polynomial}
Assume the two-radius analytic hypotheses of Theorem~\ref{III:U:thm:oracle}
and the hypotheses of Theorem~\ref{III:U:thm:normalized-transfer}, with
fixed effective margins. Suppose the section coefficients can be
enclosed to $B$ bits in time polynomial in $m,B$ and their input
description. Then, for every $p\ge1$, a terminating algorithm returns
a dyadic polynomial $\widehat h_p$, positive on $I$, together with
certified error $\|\widehat h_p-h\|_\infty\le2^{-p}$.
Its horizon, degree and coefficient precision are $O_H(p)$.
Compilation has polynomial bit cost. Horner evaluation takes
$O_H(p\mathsf M(p))$ bit operations and $O_H(p)$ scratch bits,
with $O_H(p^2)$ bits of stored polynomial coefficients; here
$\mathsf M(p)$ is the cost of multiplying $O(p)$-bit integers.
These costs concern the eigenfunction approximator, not the full state compiler.
\end{corollary}
\begin{proof}
Choose $n=O_H(p)$ so that the right side of
\eqref{III:U:eq:normalized-transfer}, multiplied by $e^{L'}$, is at most
$2^{-p-2}$. Let $a_*=2g_-^{s_+}$ be the lower growth bound of
Theorem~\ref{III:U:thm:oracle}. For a rounded polynomial iterate with
$E_n=\|\widehat P_n-V_n\|_\infty\le a_*^n/2$, normalization gives
\[
 \left\|\frac{\widehat P_n}{\widehat P_n(0)}-
                  \frac{V_n}{V_n(0)}\right\|_\infty
 \le 2(1+e^{L'})E_n/a_*^n.
\]
Indeed $V_n/V_n(0)\in K_{L'}$ normalized at zero and
$\widehat P_n(0)\ge a_*^n-E_n$.
Take $E_n\le a_*^n2^{-p-2}/[2(1+e^{L'})]$.
Equation~\eqref{III:U:eq:transfer-rounded-error}, with
$\delta_m=O_H(R^{-m})$, then permits $m=O_H(p)$ and per-step
certified arithmetic error $2^{-O_H(p)}$. The polynomial coefficient
norms grow at most exponentially in $n$ and normalization divides
by a quantity at least $a_*^n/2$, also exponential in $n$.
Thus $O_H(p)$ coefficient bits suffice after allowing for the
$O(p)$ operations per row and final dyadic rounding of the normalized
coefficients, whose error is included in the budget. This can be
implemented by interval arithmetic with directed dyadic rounding
and the stipulated coefficient enclosures; double precision until
the explicit error tests pass. The preceding bounds prove termination.

If necessary request the smaller final error
$\min(2^{-p},e^{-L'}/2)$; this changes precision only by a model
constant and proves positivity from $h\ge e^{-L'}$.
There are $O(nm^2)=O_H(p^3)$ dense scalar operations after section
construction, on $O_H(p)$-bit numbers; this and the coefficient oracle
give polynomial compilation cost. There are $O(p)$ coefficients of
$O(p)$ bits. Horner evaluation at a rational argument known to the
required $O_H(p)$ bits uses $O(p)$ multiplications; the coefficient
norm bound controls error accumulation. Its working accumulator uses
$O_H(p)$ bits. Immutable coefficients are charged separately.
\end{proof}

For the homographic weights of Section~\ref{III:U:sec:oracle}, one can use
$c=\max_a 2s_+|\xi_a|/(1-|\xi_a|)$ and a verified branch derivative
bound $\gamma<1$. Rational fixed exponents and rational positive
parameters give effective power enclosures by standard interval
root or logarithm/exponential algorithms. Alternatively their
certified coefficient oracle is an explicit input to the corollary.

\paragraph{Implications for state compilation.}
The normalized target, a positive lower bound, geometric convergence,
and its polynomial approximation are now supplied under explicit
effective hypotheses. A state compiler has additional acceptance
conditions, inverse-coordinate operations and inductive update
relations. The above theorem does not verify them. In particular its
fast Horner cost cannot be relabeled the cost of every compiled state update.

\section{Controlled recurrent states and exact residual coding}\label{III:U:sec:gdn}
A word-Hankel theorem for exogenously supplied letters is not automatically
a theorem for a trained decoder whose controls depend on its current
activations. This section isolates the valid matrix identity from that
unproved transfer. Gated DeltaNet combines a decay gate with a delta-rule
state update~\citep{yang2025gated}. The controlled example below uses
the unit-decay delta rule; it does not assume that the learned gates of
a trained GDN layer realize that special case.

\subsection{An executable core and a reduced residual}
Write a controlled recurrence in key-by-value orientation as
$S^+=AS+U$. Let $c$ be a retained core, $D(c)$ a computable matrix
predictor, $B^\top B=I_r$, and $Z=(S-D(c))B$.
\begin{proposition}[Closed projected residual update]\label{III:U:prop:core}
Under the same controls and the same supplied core trajectory,
\[
 Z^+=AZ+[AD(c)+U-D(c^+)]B,
 \qquad \widehat y=q^\top D(c)+q^\top ZB^\top.
\]
The readout error is
$q^\top(S-D(c))(I-BB^\top)$.
\end{proposition}
\begin{proof}
Substitute $S^+=AS+U$ into $(S^+-D(c^+))B$ and add and subtract
$AD(c)B$. The readout follows by projecting the residual on the value
side. No invariant-subspace assumption on $\operatorname{range}D$ is
needed for this identity.
\end{proof}

The core update must itself be computed from available deployed quantities.
It cannot silently query a deleted reference state. Further, the identity
is over exact arithmetic; matching FP32 bits requires an explicit
operation-order analysis or a reversible word-level residual.

\begin{proposition}[Weighted core-conditioned prediction]
Let $O$ be a specified frozen future-readout operator and $W=O^\top O$.
Assume the needed weighted second moments exist, and take a finite-valued
core. Define
\[
 D_*(c)=\mathbb E[W\mid c]^\dagger\mathbb E[WS\mid c],\qquad
 M_c=\mathbb E[(S-D_*(c))^\top W(S-D_*(c))].
\]
Then
\[
 \inf_{D,\,B^\top B=I_r}
 \mathbb E\|O(S-D(c))(I-BB^\top)\|_F^2
 =\sum_{j>r}\lambda_j(M_c).
\]
\end{proposition}
\begin{proof}
Conditional weighted normal equations eliminate the cross terms. For
any $D$, its residual Gram equals $M_c$ plus the positive semidefinite
matrix
$\mathbb E[(D-D_*)^\top W(D-D_*)]$. Null directions of
$\mathbb E[W\mid c]$ are also null for $\mathbb E[WS\mid c]$, so the
pseudoinverse solves the normal equations. Choosing $D_*$ and the top
$r$ eigenvectors of $M_c$ minimizes the omitted trace.
\end{proof}

This is conditional weighted regression plus principal components, not a
new general compression theorem. The operator $O$ must cover the future
queries and intervention cuts for which a guarantee is claimed. A richer
core can lower the optimal residual Gram, but its bits, predictor, basis,
and execution cost must be paid for.

\subsection{A constructive controlled example}
For key and value dimensions both $n\ge2$, take
$S=cI_n+a e_1^\top$, $D(c)=cI_n$, $B=e_1$, with a finite-valued
stored scalar $c$. This state is reachable from zero with $n$
rank-one delta writes: $k=e_i$, $\beta=1$, and $v=ce_i+a_i e_1$.
Subsequently allow a unit key $k$, $\beta\in[0,1]$, and
$v=ck+\xi e_1$. The unit-decay delta update gives
\[
 a^+=a+\beta k(\xi-k^\top a),\qquad
 y=cq^\top+(q^\top a)e_1^\top.
\]
Thus $c$ and $n$ residual coordinates suffice for exact continuation of
this controlled family without storing a dense $cI_n$ or dense basis.
If $c,a$ are independent centered variables, $\operatorname{Var}c=\sigma_c^2$,
and $\mathbb E\|a\|^2=V$, then for $W=I$ the raw value Gram is
$\sigma_c^2I+Ve_1e_1^\top$. A raw rank-$r$ projection loses
$(n-r)\sigma_c^2$ for $r\ge1$, whereas the core-conditioned residual
has rank one. This is an executable synthetic separation, not a finding
that Qwen exhibits this family.

\subsection{The missing mathematical bridge to encoded length}
An exact residual supplies injectivity. The predictor changes residual
word statistics, not exactness. A future-readout norm can ignore state
components that a lossless full-state codec must retain. Small floating
values can still contain many independent mantissa bits. A positive coding
theorem therefore needs assumptions on the encoded integer residual,
small-magnitude behavior, predictor description length, and adaptation
cost. A small Hankel singular tail alone does not provide them.

The next useful experiment is a frozen comparison of a deployable dynamic
core, an energy predictor, cross approximation, and existing exact residual codecs, with the
same residual format, complete bytes, and complete update/read cost.
Selecting on total encoded length is a new candidate that needs new held-out
data. For a trained-model claim, controls, gates, normalization, all caches,
logits, and free-generation token divergence must be covered, and the
deleted native state must actually be absent from resident memory.
\subsection{A covariance surrogate and what exact coding actually permits}
\label{III:U:sec:codec-covariance-audit}
Covariance does not determine exact expected code length. The valid
algebraic inequality is the following surrogate bound.

\begin{proposition}[Classical diagonal-covariance surrogate inequality]
\label{III:U:prop:covariance-surrogate}
For $\Gamma\succeq0$, $a,b>0$, and a real orthogonal matrix $U$,
\begin{equation}
 \sum_j\log_2\bigl(a+b(U^\top\Gamma U)_{jj}\bigr)
 \ge\log_2\det(aI+b\Gamma)
 =\sum_j\log_2(a+b\lambda_j(\Gamma)).
 \label{III:U:eq:covariance-surrogate}
\end{equation}
Equality holds exactly when $U^\top\Gamma U$ is diagonal, including
arbitrary rotations within repeated-eigenvalue eigenspaces.
This is an inequality for a covariance surrogate, not an exact
lossless-entropy formula.
\end{proposition}
\begin{proof}
Apply Hadamard's determinant inequality to the positive definite matrix
$aI+bU^\top\Gamma U$. Its determinant is invariant under the orthogonal
change of basis. Equality in Hadamard's inequality holds exactly when
the Gram vectors are orthogonal, equivalently when this matrix is
diagonal. Taking base-two logarithms proves the result.
\end{proof}


\begin{proposition}[A valid variance bound for discrete lossless symbols]
\label{III:U:prop:discrete-entropy-bound}
Let $Z=(Z_1,\ldots,Z_D)$ be an integer-valued residual vector with
finite coordinate variances $v_j$. Then
\begin{equation}
 H(Z)\le\sum_jH(Z_j)
 \le\frac12\sum_j\log_2\!\bigl(2\pi e(v_j+1/12)\bigr).
 \label{III:U:eq:discrete-entropy-bound}
\end{equation}
An ideal prefix code on the whole vector has expected length between
$H(Z)$ and $H(Z)+1$; coding coordinates separately instead admits
length less than $\sum_jH(Z_j)+D$. All model, predictor, basis,
framing, and coding-distribution descriptions are additional costs.
The same statements apply conditional on available side information,
by applying the bounds to each conditional law and averaging.
\end{proposition}
\begin{proof}
The first inequality is entropy subadditivity. For an integer variable
$Z_j$, add an independent uniform variable on $[-1/2,1/2]$.
The resulting density has differential entropy $H(Z_j)$ in bits and
variance $v_j+1/12$. The Gaussian maximum-entropy bound gives the
second inequality; this is the classical discrete variance bound
\citep{massey1988}. The coding claims are the prefix-code entropy
lower bound and the Shannon lengths $\lceil-\log_2p\rceil$.
\end{proof}


\paragraph{Counterexamples and missing codec assumptions.}
Covariance alone cannot determine discrete entropy. The symmetric
variable $Z=\pm1$ has variance one and entropy one bit. In contrast,
\[
 P(Y_m=0)=1-m^{-2},\qquad
 P(Y_m=m)=P(Y_m=-m)=\frac1{2m^2}
\]
also has variance one but entropy
$h_2(m^{-2})/\log2+m^{-2}\to0$.
More directly, the two-point variable $Z=\pm1$ remains exactly
representable on grids $\Delta_q=1/m$ with entropy one bit, whereas
$\tfrac12\log_2(1+2\pi e/\Delta_q^2)$ diverges. An
$O(\Delta_q^2)$ remainder cannot repair this claimed equality for
such a discrete source.
Thus an equality inferred solely from covariance, or an exact saving
threshold based on it, is false without further distributional and coding assumptions.

Nor does a nonidentity readout weight force a different basis. For
example, $\Gamma=\Udiag(1,2)$ and a diagonal weight
$W=\Udiag(2,3)\ne I$ share the coordinate eigenbasis with their
corresponding weighted covariance. No strict surrogate gap follows
merely from $W\ne I$. A measured saving would require empirical evidence
or additional structural assumptions.

An arbitrary orthogonal change of basis followed by coordinate
quantization is not an exact integer-lattice codec. For example, rotate
$(1,0)$ and $(2,0)$ by $45$ degrees; their coordinates are respectively
$(1/\sqrt2,1/\sqrt2)$ and $(\sqrt2,\sqrt2)$, both rounding to
$(1,1)$ on the unit grid. Exact reconstruction requires an injective
representation, such as an appropriate reversible integer transform
or an additional exact residual. A finite-precision basis must itself
be serialized and its actual decoder used.

For a residual matrix $R$ with $n_k$ rows,
$\mathbb E[R^\top R]$ aggregates row second moments. It is not the
per-row covariance used by a scalar marginal-entropy formula; identical
centered row laws would require division by $n_k$. Likewise, the
dimension $n_vr-r(r+1)/2$ of a Stiefel manifold is not the bit length
of an exact basis header. Actual encodings, including a completion
basis when needed, determine that length.

For a fixed exact codec the genuine saving condition is simply
\[
 L_{\rm metadata}+L_{\rm core}+\mathbb E L_{\rm residual}
                         <\mathbb E L_{\rm baseline},
\]
with the same scope and all shared costs amortized explicitly.
Equation~\eqref{III:U:eq:covariance-surrogate} can guide a heuristic or bound
a specified surrogate; it cannot replace these code lengths. For an
injective discrete representation $S\mapsto(C,Z)$,
$H(S)=H(C,Z)=H(C)+H(Z\mid C)$: an invertible representation does
not lower the fundamental joint lossless entropy, although it may
improve a restricted practical coder or exploit already available side
information.


\subsection{Reversible integer lifting with an actual coding contract}
\label{III:U:sec:integer-lifting}
Integer lifting is classical \citep{calderbank1998}. Its relevance here
is that it repairs injectivity of the proposed residual transform;
it does not by itself establish a Hankel-specific compression gain.

\begin{lemma}[Constructing lifting factors from a supplied orthogonal basis]
\label{III:U:lem:rotation-lifts}
A supplied computable $U\in SO(D)$ can be approximated arbitrarily
closely by a product of at most $3D(D-1)$ rational shears. A planar
rotation of angle $\theta$ is
\[
 \begin{pmatrix}1&-\tan(\theta/2)\\0&1\end{pmatrix}
 \begin{pmatrix}1&0\\\sin\theta&1\end{pmatrix}
 \begin{pmatrix}1&-\tan(\theta/2)\\0&1\end{pmatrix}.
\]
Split angles if necessary so that $|\theta|\le\pi/2$; each shear
coefficient then has magnitude at most one. If all $\ell$ coefficients
are rounded within $\delta$ while staying in $[-1,1]$, the product
error in maximum-row-sum norm is at most $\ell2^{\ell-1}\delta$.
\end{lemma}
\begin{proof}
Givens elimination expresses a special orthogonal matrix as at most
$D(D-1)/2$ planar rotations, allowing angles up to $\pi$.
Splitting each into at most two proves the shear count by the displayed
identity. Telescope the product; all factors before and after rounding
have norm at most two, and each factor error has norm at most $\delta$.
\end{proof}


\begin{theorem}[Bijective lifting and a covariance bound including correlated errors]
\label{III:U:thm:integer-lifting}
Let $S_r=I+s_re_{i_r}e_{j_r}^\top$, $i_r\ne j_r$, be rational
shears, $1\le r\le\ell$, and use a fixed deterministic nearest-integer
rule. Apply $z_{i_r}\leftarrow z_{i_r}+\operatorname{round}(s_rz_{j_r})$
in the stated order. The resulting map $T:\mathbb Z^D\to\mathbb Z^D$
is a bijection, whose inverse applies the same rounded values with
subtraction in reverse order. Put $U_*=S_\ell\cdots S_1$ and
\[
 b_j=\frac12\sum_{r=1}^\ell
       |(S_\ell\cdots S_{r+1})_{j,i_r}|.
\]
For a random integer vector $Z$ of finite second moment and covariance
$M$, the output obeys
\begin{equation}
 \operatorname{Var}(T(Z)_j)
 \le\left(\sqrt{(U_*MU_*^\top)_{jj}}+b_j\right)^2.
 \label{III:U:eq:lifting-variance}
\end{equation}
If a desired orthogonal $U$ diagonalizes $M$ with diagonal entries
$\lambda_j$, the right side is at most
\[
 \left(\sqrt{\lambda_j}+
    \|(U_*-U)_{j,:}\|_2\sqrt{\Utr M}+b_j\right)^2.
\]
These assertions apply separately to conditional laws given a core.
\end{theorem}
\begin{proof}
The donor coordinate $j_r$ is unchanged by a shear, so subtraction
recovers its integer target exactly. Reverse composition proves
bijectivity, for any deterministic rounding convention.
Writing the rounding error at step $r$ as $\epsilon_r e_{i_r}$,
$|\epsilon_r|\le1/2$, gives
\[
 T(z)=U_*z+\sum_{r=1}^\ell
            S_\ell\cdots S_{r+1}e_{i_r}\epsilon_r(z).
\]
Thus the $j$th error lies in $[-b_j,b_j]$ and its standard deviation
is at most $b_j$. The triangle inequality in centered $L^2$ proves
\eqref{III:U:eq:lifting-variance}, without independence. For the last
bound, split $U_*=U+(U_*-U)$ and apply the same inequality, using
$\operatorname{Var}(vZ)\le\|v\|_2^2\Utr M$.
\end{proof}

The claim that each lift adds at most $1/4$ to the \emph{output}
variance is false. Let $Z_1=Z_2$ be Bernoulli$(1/2)$ and lift with
$s=51/100$. The exact linear target has variance $(151/100)^2/4$,
whereas the integer target is $2Z_1$, with variance one. The increase
is $17199/40000>1/4$, because the rounding error is correlated with
the signal. The propagated-error terms above are indispensable.

\begin{theorem}[Complete expected length against a specified baseline]
\label{III:U:thm:actual-lifting-code}
Let $S$ be a finite-bit integer vector and $C=c(S)$ a core. Given $C$,
an integer predictor and a specified lifting transform produce
$Z'=T_C(S-D(C))$. The decoder receives the core and every model,
predictor, transform and framing description needed for inversion;
their total nonpayload length is at most $B_{\rm meta}$ bits.
Fix the coding models before coding the evaluated source; any dependence
on the transmitted core is explicitly conditional. Let prefix codes
for $C$ and for each $Z'_j$ conditional on $C$ have
Shannon lengths for specified full-support models $Q_C,Q_{j|C}$.
Use base-two KL and put
\[
 D_C=D_2(P_C\|Q_C),\qquad
 D_j(c)=D_2(P_{Z'_j|c}\|Q_{j|c}).
\]
If $v_j(c)$ is any valid bound on the conditional variance, the
expected complete length is at most
\begin{align}
 \mathbb E L\le B_{\rm meta}+H(C)+D_C+1+D
 +\mathbb E_C\sum_{j=1}^D\left[
 \tfrac12\log_2(2\pi e(v_j(C)+1/12))+D_j(C)\right].
 \label{III:U:eq:actual-lifting-code}
\end{align}
If this explicit upper bound is smaller than a proved lower bound
on the expected length of a \emph{specified} baseline coder, the new
coder improves that baseline in expectation at exact reconstruction.
\end{theorem}
\begin{proof}
The inverse lifting steps and the integer predictor reconstruct $S$
exactly from $C,Z'$. Expected Shannon lengths are bounded by cross
entropy plus one for each code: cross entropy is entropy plus the
displayed model mismatch. Proposition~\ref{III:U:prop:discrete-entropy-bound}
bounds each marginal conditional entropy by its variance term.
Sum these inequalities, average over the core, and add every
nonpayload description. The final comparison is then immediate.
\end{proof}

This is a coding theorem for the declared models, not an assertion
that their probabilities or description lengths are free to obtain.
An implementation must also specify finite integer widths. For
bounded source support, the lifting recurrence bounds every intermediate
magnitude by $\|S_r\|_\infty$ times the previous bound plus $1/2$;
these explicit bounds determine exact numerator, product and accumulator
widths. In-place lifting stores $D$ integers plus multiplication scratch;
holding additional input/output buffers must be charged. No reduction
of update latency or peak memory is implied by \eqref{III:U:eq:actual-lifting-code}.

\begin{proposition}[The universal saving threshold is impossible]
\label{III:U:prop:no-universal-code-gain}
For the deterministic exact representation above,
$H(C,Z')=H(S)$. Consequently a valid complete-code upper bound cannot
lie below a valid lower bound on $H(S)$. In particular a
criterion for beating \emph{every} lossless coder by placing its
complete length below an entropy lower bound has no realizable instance.
\end{proposition}
\begin{proof}
The encoder is deterministic and injective, and has a deterministic
inverse on its image; entropy is invariant. Any uniquely decodable
code has expected length at least $H(S)$. If $U$ bounds its complete
length above and $L$ bounds $H(S)$ below, then
$U\ge\mathbb E L_{\rm code}\ge H(S)\ge L$, so $U<L$ is impossible.
\end{proof}

For comparison, let $S=(Z,Z)$ with $Z$ uniform on $b$-bit integers.
The lift $(z_1,z_2)\mapsto(z_1,z_2-z_1)$ yields $(Z,0)$.
With the distribution and lift fixed, $b$ payload bits suffice instead
of a raw $2b$-bit layout; including a header of length $H_0$ still
saves bits if $H_0<b$. The joint entropy was already $b$, so this is
a gain over the specified raw layout, not over an optimal lossless code.

For a practical fixed frame, a one-bit selector between a new exact
codec and an exact baseline gives length at most
$1+\min(L_{\rm new},L_{\rm base})$, with all headers included.
This provides a deterministic fallback guarantee; it does not supply
held-out evidence for learned Hankel predictors.

\section{The homogeneous norm identity}
Here $d=|\Sigma|$ and $C_{i,j}(f)=(f(uv))_{|u|=i,|v|=j}$.
\begin{lemma}[Homogeneous norm and spectrum]\label{III:S:lem:homogeneous}
If $f\ne0$ is supported at length $m$, then $\Snorm{H_f}=a:=\Snorm f_2$. If $m\ge2$, then $\sigma_1=\sigma_2=a$ and $\sigma_3\ge a/\sqrt d$.
\end{lemma}
\begin{proof}
The word-length decomposition makes each nonzero block map degree $q$ to degree $m-q$, with block $C_{m-q,q}(f)$. Their input and output spaces are mutually orthogonal, so their singular-value lists combine. Each Frobenius norm is $a$. The two endpoint blocks have singular value $a$, and every other block has norm at most $a$. The block $C_{1,m-1}$ has at most $d$ rows and Frobenius norm $a$, hence has a singular value at least $a/\sqrt d$.
\end{proof}


\section{Conditional counterexamples and exact sanity checks}\label{III:S:app:conditional}
All examples below use completed-string semantics with $v=\mathbf1$ and initial row uniform. The matrices are nonnegative and satisfy $\beta+\sum_aA_a\mathbf1=\mathbf1$, with strictly positive termination probabilities. They therefore define proper finite-string distributions.

\paragraph{A simple Perron root without positive overlap is insufficient.}
Take
\[
 A_j=\Sdiag(1/2,1/4),\quad
 A_a=\tfrac18\mathbf1e_1^\top,\quad
 A_b=\tfrac18\mathbf1e_2^\top,\quad\beta=(1/4,1/2)^\top.
\]
Both pure normalized states are reached by reset letters. Repeated $j$ leaves them pure, with current termination probabilities differing by $1/4$ for every run length. Thus $\widehat\omega_{j,N}\ge1/4$, although the eigenvalue ratio of $A_j$ is $1/2$. The dominant right eigenvector has a zero coordinate; the second state has zero Perron overlap. This is exactly the hypothesis missing from a simple-eigenvalue-only claim.

\paragraph{A Jordan factor is necessary even with a strictly positive matrix.}
In dimension three put $P=\mathbf1\mathbf1^\top/3$, $u=(1,-1,0)^\top$, $w=(1,1,-2)^\top$, $L=uw^\top$, and
\[
 A_j=\tfrac14I+\tfrac14P+\tfrac1{48}L,\quad
 \beta=(1/8,1/4,3/8)^\top,\quad
 g=(3/16,1/8,1/16)^\top,\quad A_a=ge_1^\top,\ A_b=ge_2^\top.
\]
Here $A_j>0$, $A_j\mathbf1=\mathbf1/2$, $L^2=PL=LP=0$, and both reset states are reachable. Expanding the Jordan power gives
\[
 (e_1^\top-e_2^\top)(2A_j)^N\beta
       =-\frac{N+2}{16}\,2^{-N}.
\]
This is a coordinate of the conditional future difference. Hence no bound $K2^{-N}$ with a fixed $K$ holds, although the exact root rate is $1/2$.

\paragraph{Absolute contraction may leave conditional memory unchanged.}
For
\[
 A_j=I/2,\quad A_a=E_{11}/4,\quad A_b=E_{22}/4,\quad
 \beta=\mathbf1/4,
\]
one has $Q=\begin{psmallmatrix}1/11&1/12\\1/12&1/11\end{psmallmatrix}$ and
$\widehat\omega_N^2=1/66$ for all $N$. The window $j^N$ retains the two pure normalized states, and the simplex has no larger diameter. The raw transition JSR is $1/2$. This example independently checks that the metric used in Section~\ref{III:S:sec:metric} cannot simply be relabeled as a conditional one.
